\documentclass[paper=b5,fontsize=11pt,DIV=11,open=any,headinclude=true,numbers=noenddot]{scrbook}
\usepackage{snegmanual}
\addbibresource{references.bib}
\hypersetup{pdftitle={Second Quantization with SNEG},pdfauthor={SNEG documentation project}}
\begin{document}
\frontmatter
\title{Second Quantization with SNEG}
\subtitle{A practical guide for solid-state theory}
\author{SNEG documentation project}
\date{Edition for SNEG 2.2}
\publishers{\small Licensed under \href{https://creativecommons.org/licenses/by/4.0/}{Creative Commons Attribution 4.0 International}.}
\maketitle
\addchap{Preface}
An electron hopping between two orbitals is described by a bilinear.
The same bilinear becomes a current when differentiated with respect to
a bond phase, and an exchange interaction when iterated through a
high-energy charge state. Keeping the operators until the last step
makes these connections visible. This book uses the Wolfram Language
package SNEG~\cite{sneg} to do the operator algebra and small-system
calculations that connect those descriptions.

Part I follows the construction of a calculation from declarations to
matrix elements. Part II can be read by subject: each worked calculation
specifies its Hamiltonian, units and basis. Examples headed as case
studies start an independent kernel session with SNEG loaded; evaluate
the inputs within each study in sequence. Input labels restart accordingly.
A final \fn{TableForm} or \fn{MatrixForm} formats an answer for display;
store the underlying list or matrix before applying it. Part III gives
the calling conventions of the documented symbols, with a page index
at the back. Exercises and their solutions use the same conventions as
the chapters. Energies include additive constants unless stated otherwise;
$\hbar=k_{\mathrm B}=1$.
\setcounter{tocdepth}{1}
\tableofcontents
\mainmatter
\part{Tutorial}
\chapter{Getting started}
\setwlchapter{getting-started}
If two electrons of opposite spin occupy a localized orbital,
removing the up electron must leave a down electron. Changing the
order in which the two creators were applied changes the sign of the
resulting ket, but cannot change its occupation. This distinction
between an amplitude and a probability is a good first test of an
operator calculation.

\section{Installation and a first session}
Install the release paclet with \fn{PacletInstall} in Mathematica
12.1 or later. The English help pages require Mathematica 14.0 or later.
The installation command, which contacts the release server, is:
\begin{lstlisting}[style=wl]
PacletInstall[
 "https://github.com/rokzitko/sneg/releases/latest/download/SNEG.paclet"]
\end{lstlisting}
Load the package \emph{before} entering a definition using its symbols.
The lower-case loader also works on existing SNEG notebooks. In the
Documentation Center open \texttt{paclet:SNEG/guide/SNEG}.
\begin{wl}
Needs["sneg`"];
$SnegVersion
\end{wl}
Choose a new operator name and declare its algebra. In the following
input \fn{d[]} means one spin-degenerate orbital, not an annihilator.
The latter needs both a type and a spin index.
\begin{wl}
snegfermionoperators[d];
u = d[AN, UP]; v = d[AN, DO];
TableForm[{{antikomutator[u, conj[u]], antikomutator[u, conj[v]]},
           {antikomutator[v, conj[u]], antikomutator[v, conj[v]]}},
  TableHeadings -> {{"up", "down"}, {"up", "down"}}]
\end{wl}
The unit diagonal expresses Pauli exclusion for a single spin-orbital;
the off-diagonal zeros express orthogonality. The rule says nothing
about the occupation of the orbital. That information belongs to a
state or density matrix.

\section{From an operator to a spectrum}
Give the level an energy $\epsilon$ and assign an extra cost $U$ when
both spins are present. The algebra itself is independent of their
values. A separate line declares the parameters real so that
conjugation later has the intended meaning.
\begin{wl}
snegrealconstants[\[Epsilon], U];
h = \[Epsilon] number[d[]] + U hubbard[d[]];
blocks = makematricesbzvc[h, qsbasisvc[{d[]}]];
TableForm[{#[[1, 1]], #[[1, 2]], #[[2, 1, 1]]} & /@ blocks,
  TableHeadings -> {None, {"Q", "Spin", "Energy"}}]
\end{wl}
There are four states but only three rows. The middle row represents
both projections of a spin $1/2$. The label $Q=N-1$ counts electrons
relative to one electron per site. If $U>0$, changing $\epsilon$ from
positive to negative first makes the spin doublet the ground state.
Only below $\epsilon=-U$ is it favourable to add the second electron.
Chapter~10 revisits these two addition thresholds as poles of a
Green's function.

\section{Reading operator output}
A dagger distinguishes creation from annihilation in the formatted
output. To inspect the actual expression, request \fn{InputForm};
use \fn{FullForm} to see every head and argument. Both are views.
For instance, the occupation polynomial represented by
\fn{hubbard[d[]]} is unchanged if its displayed product acquires a
minus sign from putting all creators on the left. Use \fn{snegtex}
when the result is needed as a formula in a paper. A useful final
step before export is to compare a matrix element, rather than a
particular ordering of symbols.
\casestudy{moment}{Local moment as an occupation measurement}
\begin{exercise}[Parity spectroscopy]\label{ex:first}
At $\epsilon=-U/3$, $U>0$, determine the addition and removal costs
from the ground multiplet. Which excited level would be invisible in
a calculation restricted to the same charge sector?
\end{exercise}

\chapter{Operators and their algebra}
\setwlchapter{algebra}
\section{What a declaration fixes}
The algebra of an operator is specified by its head, not its printed
letter. For fermions a mode is the tuple of indices \emph{after}
the type \fn{CR} or \fn{AN}. Changing any member of this tuple gives
an orthogonal mode. For spinful electrons the last index denotes spin;
an orbital, site, band or channel index belongs before it. Thus
$c^\dagger_{2a\uparrow}c_{2b\uparrow}$ transfers charge between
orbitals at fixed site and spin. SNEG keeps the same mode ordering
under conjugation and reverses the order of factors in a product.
\casestudy{spin-three-halves}{Changing the local spin multiplicity}

\section{Products, scalars and Hermitian conjugation}
With $\{c_i,c_j^\dagger\}=\delta_{ij}$, transposing neighbouring
factors generates a lower-degree term as well as a sign. Ordinary
Wolfram Language multiplication sorts its arguments and cannot encode
this operation. Use \fn{nc[a,b,c]} and \fn{pow[a,n]} for operator
products and their powers; use ordinary multiplication for commuting
coefficients. Declare real constants with \fn{snegrealconstants} and
complex amplitudes with \fn{snegcomplexconstants}. If a hopping
amplitude can carry flux, its conjugate is a different coefficient.
Real functions and real-valued indices are separate declarations:
$\epsilon(k)$ may be real even if its functional form is unknown.
\begin{wl}
snegfermionoperators[c];
snegcomplexconstants[v];
a = v nc[c[CR, 1, UP], c[AN, 2, UP]];
h = a + conj[a];
Expand[{h - conj[h], komutator[h, number[c[1]] + number[c[2]]]}]
\end{wl}
The two zeros answer different questions: the reverse hop is present
with the conjugate amplitude, and every term removes as many particles
as it creates. Hermitian pairing would pass the first test but generally
fail the second.

\section{Reference states}
In \fn{EMPTY} order an annihilator kills the vacuum. In \fn{SEA} order
the filled modes have negative energy or momentum index relative to
the chosen Fermi point. A creator below the surface therefore kills
the sea. At the Fermi point one must specify the state or resolve the
degeneracy; changing the order alone does not determine its physical
occupation. \fn{NONE} preserves products within a family and is useful
when explicitly comparing before and after a contraction. Changing
the ordering cannot undo simplifications of an already evaluated product.
\casestudy{sea}{A particle and a hole above the reference sea}

\section{Bosons, Majoranas and spin operators}
For a boson $[b,b^\dagger]=1$ replaces the fermion anticommutator.
A self-adjoint SNEG Majorana obeys $\{\gamma_i,\gamma_j\}=\delta_{ij}$:
its square is $1/2$, rather than one. A spin of size $s$, declared with
\fn{snegspinoperators}, acts in dimension $2s+1$ and commutes with
operators at other sites. Local products can be reduced using
$S^+S^-=s(s+1)-(S^z)^2+S^z$; a spin raising operator is not a
canonical boson. A Grassmann coefficient is yet another object:
it anticommutes with odd fermionic quantities but supplies no
creation--annihilation contraction.
\casestudy{pair-algebra}{Does a local fermion pair behave as a boson?}
\section{Transitions between many-body states}
Hubbard operators obey $X^{ab}X^{cd}=\delta_{bc}X^{ad}$ on the
same site. On different sites the parity of the transition matters:
moving an electron from one site to another still carries a Fermi
sign even if individual orbital operators have been eliminated.
\casestudy{graded}{Grading local transitions}
\begin{exercise}[Majorana normalization]\label{ex:majorana}
Write $f=(\gamma_1+i\gamma_2)/\sqrt2$. Determine $f^2$,
$\{f,f^\dagger\}$ and $f^\dagger f$ using the SNEG normalization.
\end{exercise}

\chapter{Building Hamiltonians}
\setwlchapter{models}
\section{One bond at a time}
For a real amplitude $t$, \fn{-t hop[c[i],c[j]]} includes both
directions and both spins. The first term in a sum over distinct
unordered bonds should therefore carry $-t$, not $-t/2$. On a finite
open chain the last bond ends at $L$; on a periodic chain it joins
$L$ to 1. The distinction is measurable already in the one-electron
spectrum.
\casestudy{hopping}{Standing waves and open ends}
The open-chain wave function has nodes just outside the lattice at
sites 0 and $L+1$. Its momentum denominator is consequently $L+1$.
Periodic boundary conditions instead identify sites separated by $L$.
For complex hopping write $v c_i^\dagger c_j+\hc$ explicitly or use
\fn{matrixhop}; multiplying the entire Hermitian bond by complex $v$
does not insert a Peierls phase.
\casestudy{continuity}{Testing a bond by its current}

\section{Local interactions and constraints}
The number operator, the double-occupation projector and the spin
Casimir give three distinct probes of a local shell. On one spinful
orbital $D=n_\uparrow n_\downarrow$ is a projector, not $n^2$:
$n^2=n+2D$. A penalty $\lambda(n-1)^2$ lifts both even charge states,
whereas $UD$ acts only on the doublon. Only at fixed total charge
can one freely trade a uniform term proportional to $n$ for a constant.
\casestudy{projectors}{When does the orbital carry spin?}

\section{Pairing and the Nambu representation}
A source for a singlet pair is invariant under spin rotations but not
under a change of electronic phase with its coefficient held fixed.
The Hermitian conjugate of a pair \emph{reverses} the order of the
two annihilators. A sign error at this step can turn real eigenvalues
into imaginary ones rather than simply changing the sign of the gap.
\casestudy{pairing}{A complex pairing amplitude}
In a doubled basis $\Psi=(a,a^\dagger)^T$ one writes
$H=\frac12\Psi^\dagger M\Psi+C$. The factor $1/2$ compensates for
doubling; the constant resulting from commuting $a$ through
$a^\dagger$ must be kept when comparing energies. For a short open
Hubbard chain \fn{hamiltonian} is a convenient model constructor;
its hopping parameter enters with a plus sign. The explicit bond sum
is preferable for a ring or spatially varying coefficients.
\begin{exercise}[Counting bonds]\label{ex:bond}
Compare a two-site molecule with a periodic chain of length two whose
unit-cell sum contains a forward hop and its conjugate at each site.
Find both one-particle spectra for real $t>0$. Why is dividing by
two not a general remedy for rings with $L>2$?
\end{exercise}

\chapter{Symbolic sums and momentum space}
\setwlchapter{momentum}
\section{An index belongs to its sum}
In $\sum_k a_k^\dagger a_k$ the value of $k$ has no meaning outside
the sum. Multiplying two copies requires two independent summation
indices, even if the same letter was used to write both. Restricting
them to be equal would throw away every term coupling two different
levels. The symbolic head \fn{sum} preserves this distinction;
\fn{Sum} with numerical limits instead enumerates a cluster.
\casestudy{dummies}{Nested summations}

\section{Momentum transfer and Fourier convention}
For $N_s$ Bravais cells take
\[
 c_R=N_s^{-1/2}\sum_k e^{ik\cdot R}c_k,\qquad
 c_k=N_s^{-1/2}\sum_R e^{-ik\cdot R}c_R .
\]
One independent sum supplies one factor $N_s$ when a lattice phase
is constant. A density bilinear brings $1/N_s$ but its site sum brings
$N_s$ back; a local quartic vertex retains $1/N_s$ and three independent
momenta. Declare lattice and momentum indices with the \emph{same}
lattice size. Keep \fn{blochphase} symbolic when combining deltas;
convert to exponentials only for an explicit dispersion or geometry.
\casestudy{transfer}{Following the transfer through a density operator}

\section{Periodic clusters and boundary twists}
In one dimension $c_{n+L}=e^{i\theta}c_n$ selects the grid
$k_m=(2\pi m+\theta)/L$. Bilinears involve differences of two
momenta: the twists cancel in a transferred momentum. Thus it is
incorrect to shift \emph{every} momentum in a momentum-conservation
delta by $\theta/L$. \fn{clusterhop} inserts the phase only when a
real-space hop crosses the boundary; \fn{fourierrules} gives the
explicit finite-dimensional change of basis. In a periodic untwisted
cluster \fn{sumCluster} folds indices modulo each lattice dimension.
\casestudy{twist}{Flux sensitivity of the spectrum}
\begin{exercise}[Discrete unitarity]\label{ex:grid}
Show that a common twist does not change the inner product of two
columns of the finite Fourier matrix. Test $L=3$, $\theta=\pi/2$ and
explain why a reversed sign in just the inverse transform fails.
\end{exercise}

\chapter{States, bases and exact diagonalization}
\setwlchapter{bases}
\section{Fix the mode order first}
Let $|n_1\ldots n_M\rangle=(a^\dagger_1)^{n_1}\cdots
(a^\dagger_M)^{n_M}|0\rangle$. An annihilator acting on slot $j$
picks up $(-1)^{\sum_{i<j}n_i}$ before it reaches that slot. Neither
the binary occupation vector nor the site labels alone determine this
sign: the associated \fn{BASIS} is part of the representation. Keep
the mode order when moving a matrix into another numerical code.
An operator-form state instead retains the creators explicitly and
can be inspected without decoding bit positions.
\begin{wl}
snegfermionoperators[d];
makebasis[{d[]}];
{BASIS, ap[d[AN, DO], vc[1, 1]]}
\end{wl}

\section{Block labels versus degeneracies}
Use charge and spin projection if $[H,N]=[H,S^z]=0$. Use total spin
only if $H$ commutes with all components. A transverse field can break
$S^z$; a pairing term breaks $N$ but may preserve spin and fermion
parity. Projecting onto a non-invariant block loses couplings rather
than proving that they vanish. One diagnostic for a projector $P$ is
$(1-P)HP=0$. A second is to compare the reconstructed full spectrum
for a cluster small enough that symmetries are unnecessary.
\casestudy{multiplets}{Counting states in a highest-weight basis}
\casestudy{reflection}{Separating mirror-even and mirror-odd states}
Reflection exchanges sites, not just occupation numbers. A permutation
of creators must be reordered into the reference order. This can supply
a minus sign even when the \emph{pattern} of occupied sites is invariant.
Likewise particle--hole conjugation maps $Q$ to $-Q$ and interchanges
local double occupation and the vacuum. On a bipartite lattice at
particle--hole symmetry one can use a continuous charge isospin
instead of just pairing states of opposite charge.
\casestudy{isospin}{Spin and charge multiplets on a bipartite dimer}

\section{Matrix elements, eigenstates and thermal traces}
In each invariant sector the matrix element is $\langle i|H|j\rangle$,
with the ket indexed by the column. Wolfram Language
\fn{Eigensystem} returns eigenvectors as a list of vectors. For a
numerical normalized eigenvector $v$ compute an expectation as
$v^\dagger A v$, not $v^T A v$. For an operator changing quantum
numbers, such as $c^\dagger$, a matrix with both indices in the same
sector may be identically zero. Use \fn{makeallmatricesbzvc} to obtain
the rectangular blocks between sectors.
\casestudy{thermal}{A partition function from reduced spin blocks}
\begin{exercise}[Weight of a multiplet]\label{ex:multiplet}
The two-spin Heisenberg dimer has one singlet at $-3J/4$ and one
triplet at $J/4$. Derive $\langle\vS_1\cdot\vS_2\rangle_T$ and
check both temperature limits for $J>0$.
\end{exercise}

\chapter{Expectation values and Wick's theorem}
\setwlchapter{wick}
\section{An average needs a state}
Normal ordering is an algebraic operation, whereas contraction depends
on the reference state. For a number-conserving Slater determinant the
only nonzero two-point functions in its natural-orbital basis are
$\langle c_i^\dagger c_j\rangle=\delta_{ij}n_i$ and
$\langle c_i c_j^\dagger\rangle=\delta_{ij}(1-n_i)$. In a paired
state the anomalous contractions need not vanish. Specifying the
occupations alone is then insufficient. In an interacting thermal
state even the complete set of two-point functions need not determine
a four-point function.

\section{Direct and exchange contractions}
For $c_i^\dagger c_j^\dagger c_l c_k$ in a determinant the contraction
of $i$ with $k$ and $j$ with $l$ has positive sign. Pairing $i$ with
$l$ and $j$ with $k$ requires one extra exchange. For $i=j$ these
terms cancel exactly. The same minus sign makes the exchange energy
negative for a repulsive density interaction. The case study keeps
the spin labels distinct until the contraction has been taken.
\casestudy{exchange}{The exchange energy of a filled sea}
In an oscillator vacuum each interchange contributes a plus sign.
The number of pairings now appears explicitly in higher moments.
\casestudy{oscillator}{Gaussian moments of a displacement}

\section{Finite-temperature and anomalous averages}
For a quadratic fermion Hamiltonian the exact thermal state remains
Gaussian. Diagonalizing its Nambu matrix gives a covariance matrix;
Wick contractions with that matrix recover the expectation of any
polynomial in the fermions. This route avoids a trace over $2^M$
many-body states for $M$ spin-orbitals. A number projector on the full
state generally destroys this simple Gaussian structure.
\casestudy{gaussian}{Checking the anomalous contraction}
\begin{exercise}[Wick versus a fixed-charge state]\label{ex:wick}
Compare an equal mixture of the two one-electron states of an atom
with a free spin-degenerate level at half occupation per spin.
Show that all normal one-body averages coincide but the two double
occupations differ. What information has been lost?
\end{exercise}

\chapter{Working efficiently}
\setwlchapter{workflow}
\section{Expand at the step that needs it}
For a four-fermion interaction written as $U D$ the compact notation
hides the number of terms obtained after a transformation of orbitals.
Replacing every creator by a sum of $M$ creators can produce $M^4$
terms before momentum selection. Apply Kronecker deltas while the
indices to be eliminated are still visible. A good intermediate
observable is often a commutator or a scalar matrix element rather
than the expanded Hamiltonian itself.
\casestudy{collect}{Separating expansion from collection of sums}
\casestudy{recognize}{Recovering the physical spin coupling}

\section{How large is the matrix?}
For $L$ spinful orbitals the unrestricted basis has $4^L$ states, so
a dense matrix contains $16^L$ entries. At $L=8$ a dense complex
double-precision array already requires 64 GiB just for its elements.
A sparse hopping Hamiltonian instead connects each Fock state to at
most a number of others proportional to the number of bonds. Symmetry
reduction and sparse storage address different parts of this cost.
For a spectral line one often needs only the initial sector and the
sectors connected by the probe, rather than all transitions.

\section{Tests before a long run}
Start with one bond. Check Hermiticity, a conserved quantity and
one matrix element whose sign is known. Next check a limiting value:
does the hopping spectrum reduce to on-site levels at $t=0$? Only
then increase the cluster. For degenerate eigenvalues compare spectral
projectors or subspace weights; two diagonalizers may choose different
orthonormal vectors inside the same eigenspace.

A short export example checks that the energy symbol and spin label
survive conversion to LaTeX:
\begin{wl}
snegfermionoperators[c];
snegrealconstants[\[Epsilon]];
tex = snegtex[\[Epsilon] number[c[], UP]];
StringContainsQ[tex, "\\epsilon"] && StringContainsQ[tex, "\\uparrow"]
\end{wl}
Save the mode list with any exported numerical matrix and record
whether additive constants and chemical-potential terms are included.
A sorted eigenvalue list alone cannot detect a reversed site order.
\begin{exercise}[An independent matrix test]\label{ex:trace}
Find the trace of $U\sum_i D_i$ over the unrestricted Fock space of
$L$ spinful orbitals. Would checking only that trace detect an
incorrect sign in the hopping? Give a second trace that would
distinguish missing bonds.
\end{exercise}

\part{Applications in solid-state theory}
\chapter{Electrons in crystals}
\setwlchapter{crystals}
\section{From the gas to a lattice}
For the three-dimensional unpolarized electron gas let
$4\pi(r_s a_0)^3/3=1/n$, so that $k_F=(3\pi^2 n)^{1/3}$. The Fermi
sphere gives the kinetic energy per electron $3\hbar^2 k_F^2/(10m)$.
To see the different density scaling of exchange, integrate
$4\pi e_*^2/q^2$ over two occupied spheres of equal spin, where
$e_*^2=e^2/(4\pi\epsilon_0)$. At fixed separation $q$ the overlap
volume is $\pi(4k_F+q)(2k_F-q)^2/12$ for $q\le2k_F$ and zero otherwise.
The radial measure cancels the interaction denominator. Dividing by
the density and including the exchange factor $-1/2$ gives
$E_x/N=-3e_*^2 k_F/(4\pi)$. The neutralizing background removes the
Hartree energy. Thus kinetic and exchange terms scale respectively as
$r_s^{-2}$ and $r_s^{-1}$; their competition alone does not give the
exact correlation energy.

On a Bravais lattice with hopping $-t$ on each nearest-neighbour bond,
$\epsilon_k=-2t\sum_{\alpha=1}^{d}\cos k_\alpha$ for the chain,
square and simple cubic lattices with unit spacing. Bandwidth $4dt$
does not determine the density of states: at a stationary point its
shape also depends on the signs of the local curvatures. A 2D saddle
gives the logarithmic van Hove singularity at zero energy on the square
lattice; the 1D band edge instead has an inverse-square-root divergence.
\casestudy{square}{A square-lattice Bloch matrix}

\section{Orbital basis, valleys and topology}
Take one $A$ and one $B$ orbital per cell. In the SSH chain the
off-diagonal matrix element in a cell-periodic convention is
$q(k)=v+w e^{-ik}$. Its path is a circle about $v$. With chiral
symmetry and a fixed cell and termination, this path encloses the
origin if $|w|>|v|$; one cannot change that winding without crossing
$q=0$. The spectrum $\pm|q|$ is insensitive to the orientation of
the path, but the phase of an eigenvector is not. At $v=w$ the gap
closes at $k=\pi$ for real positive hoppings.
\casestudy{ssh}{Where the dimerized-chain gap closes}
\subsection{Measuring a winding without choosing eigenvector phases}
An eigensolver may multiply its eigenvector by a different phase at
every momentum. Direct finite differences of these vectors are then
meaningless. In a two-band chiral model with no diagonal terms use
the phase of $q(k)=h_{AB}(k)$ instead. Along a closed mesh add
$\arg[q(k_{j+1})/q(k_j)]$, taking each step on the principal branch.
The unit cell and ordered pair of orbitals must stay fixed. For the
orientation $q=v+w e^{-ik}$ the nontrivial integer is negative.
The occupied-band Berry phase modulo $2\pi$ is $\pi$ times this
winding modulo $2\pi$ in the cell-periodic convention.
\begin{wl}
winding[v_,w_,n_] := Module[{q = N[v + w Exp[-2 Pi I Range[0,n-1]/n]]},
  Round[Total[Arg[RotateLeft[q]/q]]/(2 Pi)]];
TableForm[Table[{v,1,winding[v,1,40],winding[v,1,80]}, {v,{1/2,2}}],
 TableHeadings -> {None,{"v","w","40 points","80 points"}}]
\end{wl}
At $v=w$ the mesh includes $q(\pi)=0$, so the winding is undefined;
do not repair division by zero by adding a small imaginary shift.
A phase step can cross the branch cut on a coarse mesh even short of
a closing. Refining the mesh or following the circle analytically
distinguishes this sampling error from a transition.

\subsection{A second translation vector}
In graphene the two independent translation vectors supply a second
phase. At a valley the three complex hopping amplitudes sum to zero,
and linearizing their sum gives a Dirac cone. A staggered sublattice
potential gives equal mass signs in the two valleys for a fixed spin;
intrinsic Kane--Mele spin-orbit coupling gives opposite signs and
reverses them for the other spin. Time reversal therefore allows
opposite spin-resolved Chern numbers with zero total Chern number in
the spin-conserving limit.
\casestudy{graphene}{Destructive interference at a Dirac point}
\casestudy{rashba}{Which spin component splits the band?}

\section{Flux per plaquette}
In Landau gauge the forward $y$ hop at coordinate $x$ carries
$e^{2\pi i\alpha x}$ for flux $\alpha$ in units of a flux quantum per
plaquette. For $\alpha=p/q$ in lowest terms a magnetic cell contains
$q$ sites. At fixed $k_y$ the diagonal entries of the Harper matrix
are $-2t\cos(k_y+2\pi p x/q)$; nearest $x$ neighbours are coupled
by $-t$ with a boundary twist $e^{iqk_x}$. Here energies are in units
of $t$; at the centre of the magnetic Brillouin zone:
\begin{wl}
harper[q_, ky_] := Table[-2 Cos[ky + 2 Pi i/q] KroneckerDelta[i, j] -
  Boole[Mod[i-j, q] == 1] - Boole[Mod[j-i, q] == 1],
  {i, 0, q-1}, {j, 0, q-1}];
harper[3, 0] // MatrixForm
\end{wl}
The oriented product of link phases around a plaquette is gauge
invariant; individual bond phases need not be. The same distinction
matters for current operators derived by differentiating the hopping
with respect to vector potential~\cite{vucicevic}.
\begin{exercise}[A magnetic cell]\label{ex:harper}
For flux $1/3$ find the three energies at $k_x=k_y=0$. Check their
sum and the sum of their squares without using the eigenvectors.
\end{exercise}

\chapter{Lattice vibrations and electron--phonon coupling}
\setwlchapter{phonons}
\section{Quantizing a displacement}
For a monatomic chain with spring constant $K$ and mass $M$,
$\omega_q=2\sqrt{K/M}\,|\sin(qa/2)|$. For $\omega_q>0$ the normal
coordinate is proportional to $(b_q+b^\dagger_{-q})/\sqrt{2M\omega_q}$;
the adjoint momentum involves their difference. Pairing $q$ with $-q$
keeps the real-space displacement Hermitian. The centre-of-mass
translation at $q=0$ is a free coordinate, not an oscillator for which
one may substitute $\omega=0$ into this normalization.

\section{Density and bond couplings}
A local Holstein term $g(n_i-n_0)(b_i+b_i^\dagger)$ describes a change
of site energy. On a bond, modulation of the hopping instead gives
$\alpha(u_{i+1}-u_i)\sum_\sigma(c^\dagger_{i\sigma}c_{i+1,\sigma}+\hc)$.
Uniform displacement drops out of the bond coupling. Fourier transforming
it supplies a factor $e^{iqa}-1$ that a local optical-mode vertex does
not have. In either case a phonon carries the momentum transferred
between electron states.
\casestudy{holstein}{Checking the momentum of the emitted phonon}

\section{Displaced oscillator and induced attraction}
For a fixed electronic occupation put $B=b+g(n-n_0)/\omega_0$. Then
\[
 \omega_0 b^\dagger b+g(n-n_0)(b+b^\dagger)
 =\omega_0 B^\dagger B-g^2(n-n_0)^2/\omega_0 .
\]
The ground energy depends on charge even when the bare phonon has no
dispersion. The second difference between energies of charge 2, 1, 0
receives a shift $-2g^2/\omega_0$. In a hopping system a change of
occupation also changes the displacement; the transformed hopping
carries displacement operators. Replacing them by their phonon-vacuum
expectations is an approximation. In an atomic charge sector the
completion of the square is exact.
\casestudy{lang-firsov}{Convergence with the number of phonon quanta}
The displaced vacuum contains a Poisson distribution in the original
oscillator basis with mean $[g(n-n_0)/\omega_0]^2$. The largest
occupied number in a truncated eigenvector is a useful convergence
diagnostic in addition to its energy. On a dimer check the convergence
of the bonding--antibonding splitting itself: it can be much smaller
than the shift in an individual energy.

\subsection{How a phonon cloud changes the hopping}
Take one spinless electron and two local Einstein oscillators on a
symmetric dimer. In the absence of tunnelling the two localized
ground states are degenerate. Moving the electron interchanges which
oscillator is displaced. Their phonon wave functions have overlap
$e^{-(g/\omega_0)^2}$: each of two oscillators contributes
$e^{-(g/\omega_0)^2/2}$. Degenerate perturbation theory gives a
splitting $2t e^{-(g/\omega_0)^2}$ to first order in $t$, with coupling
$g n_j(b_j+b_j^\dagger)$ and real positive hopping. This splitting is
a different observable from the single-site relaxation energy.

Build the matrix directly as a check on the phonon convention. The
tensor order is electron site, phonon on the left, phonon on the
right; $t$ and $g$ are in units of $\omega_0$.
\begin{wl}
split[n_,t_,g_] := Module[{id,nop,x,ph,h},
  id = IdentityMatrix[n+1]; nop = DiagonalMatrix[Range[0,n]];
  x = Normal[SparseArray[{Band[{1,2}] -> Sqrt[Range[n]],
                  Band[{2,1}] -> Sqrt[Range[n]]},{n+1,n+1}]];
  ph = KroneckerProduct[nop,id] + KroneckerProduct[id,nop];
  h = KroneckerProduct[IdentityMatrix[2],ph] -
      t KroneckerProduct[PauliX,id,id] +
      g (KroneckerProduct[DiagonalMatrix[{1,0}],x,id] +
         KroneckerProduct[DiagonalMatrix[{0,1}],id,x]);
  Differences[Take[Sort[Eigenvalues[N[h]]],2]][[1]]];
TableForm[Table[{n,split[n,0.01,0.6]/0.02, Exp[-0.6^2]}, {n,{2,4,6}}],
 TableHeadings -> {None,{"Cutoff","Splitting / 2t","Overlap"}}]
\end{wl}
Converging the cutoff at fixed $t$ approaches the finite-hopping
answer. Reducing $t$ at fixed converged cutoff is the separate test
of first-order perturbation theory.

\section{Virtual phonons and coupled modes}
At imaginary transferred frequency $\nu$, eliminating a harmonic
phonon coupled linearly to density yields
$V(q,i\nu)=-2|g_q|^2\omega_q/(\nu^2+\omega_q^2)$ for the dimensionless
coordinate $b+b^\dagger$. Only for $|\nu|\ll\omega_q$ can this be
replaced by a static attraction. A canonical transformation in powers
of $g$ must also avoid resonances between real electronic transitions
and a phonon; there the eliminated state belongs in the active space.
\casestudy{squeezing}{Stability of two coupled harmonic modes}

For an orbital doublet coupled to two degenerate distortions the
linear Jahn--Teller Hamiltonian at fixed coordinates contains
$g(Q_1\tau_z+Q_2\tau_x)$. Its lower adiabatic surface is
$K(Q_1^2+Q_2^2)/2-|g|\sqrt{Q_1^2+Q_2^2}$. The continuum of classical
minima does not imply a static symmetry-breaking expectation value in
a finite quantum molecule: the kinetic energy of the nuclei mixes
distortions. A finite two-mode Fock calculation can retain both.
\begin{exercise}[Relaxation energy]\label{ex:jt}
Find the Jahn--Teller minimum radius and lowering of the energy.
Relate the latter to the displaced-oscillator result for a single
distortion quantized with $Q=(b+b^\dagger)/\sqrt{2M\omega}$.
\end{exercise}

\chapter{Correlated electrons}
\setwlchapter{correlated-electrons}

What does it cost to move one electron between two hydrogen-like atoms?
At large separation and with one electron on each atom, the final state
contains an empty orbital and a doubly occupied one. Creating this pair
costs the local repulsion $U$. Yet a virtual excursion into precisely the
same state lowers the energy of a spin singlet. A charge excitation at
energy $U$ and a spin excitation at energy of order $t^2/U$ thus have
a common origin. The two-site Hubbard model lets us follow both without
making an approximation. We will use its exact spectrum to judge two
quite different approximations: eliminating charge fluctuations, and
replacing the interaction by a self-consistent field.

The examples form one session. Start a fresh kernel with the SNEG paclet
installed and evaluate the inputs in order. The convention for a level
is \fn{c[i]}; a single electron operator carries two additional indices,
as in \fn{c[CR,i,UP]}. Spin projections are in units of $\hbar$, and
$k_{\mathrm B}=1$.

\section{One orbital: charge, spin and particle--hole symmetry}
For one Wannier orbital on each of $L$ sites we use
\begin{equation}
 H=-t\sum_{\langle ij\rangle,\sigma}
 (c^\dagger_{i\sigma}c_{j\sigma}+c^\dagger_{j\sigma}c_{i\sigma})
 +U\sum_i D_i,\qquad D_i=n_{i\uparrow}n_{i\downarrow}.
 \label{eq:hubbard}
\end{equation}
Here $n_{i\sigma}=c^\dagger_{i\sigma}c_{i\sigma}$,
$\sigma=\uparrow,\downarrow$, and a bond occurs once in the sum.
Throughout this chapter $t$ is real and $U>0$. In a crystal $U$ is the
interaction screened by degrees of freedom excluded from this orbital
model; it need not equal an isolated-atom Coulomb integral.
Our starting point is the atomic-limit construction in
Ref.~\cite{hubbard}.

\begin{wl}
Needs["sneg`"];
snegfermionoperators[c, d];
snegrealconstants[t, U, \[Epsilon]];
atom = \[Epsilon] number[d[]] + U hubbard[d[]];
{number[d[]], hubbard[d[]]}
\end{wl}

The minus sign that may appear when $D_i$ is printed in normal order is
not an attractive interaction. Moving an annihilator through a creator
of the other spin costs one interchange. A sign-independent way to read
the result is to act on states. The following basis is labelled by
$Q=N-1$ and total spin $S$. Each sector of this atom contains a single
multiplet; its energy is the sole entry of the block.
\begin{wl}
ba = qsbasisvc[{d[]}];
rows = {Sequence @@ First[#], #[[2, 1, 1]]} & /@ makematricesbzvc[atom, ba];
TableForm[rows, TableHeadings -> {None, {"Q", "S", "Energy"}}]
\end{wl}

The empty state has energy zero, the two spin states have energy
$\epsilon$, and $\ket{2}=d^\dagger_\uparrow d^\dagger_\downarrow\ket{0}$
has energy $2\epsilon+U$. \fn{qsbasisvc} stores one representative of
each spin multiplet: the dimension of a block counts copies of the
representation, not its $2S+1$ magnetic substates.
With $\epsilon$ measured relative to the chemical potential, the atom
has a spin-doublet ground state for $-U<\epsilon<0$. At either
endpoint two charge sectors cross; differentiating a single eigenvalue
through such a crossing would give the wrong ground-state occupation.

At half filling a particle--hole transformation on a bipartite lattice
takes $c_{i\sigma}\mapsto\eta_i c^\dagger_{i\sigma}$, with $\eta_i=+1$
on one sublattice and $-1$ on the other. Hopping is invariant because
neighbours have opposite $\eta$ and distinct fermions anticommute.
Locally $n_i\mapsto2-n_i$ and $D_i\mapsto1-n_i+D_i$. Thus the
grand Hamiltonian with $\epsilon=-U/2$ is invariant. Equivalently,
\[
 U(n_{i\uparrow}-\tfrac12)(n_{i\downarrow}-\tfrac12)
 =U D_i-\tfrac U2 n_i+\tfrac U4 .
\]
Dropping the last term changes the energy origin; dropping the term
linear in $n_i$ also changes the chemical potential. Neither affects
splittings at fixed $N$, but only a constant can be ignored when
comparing absolute addition and removal energies at a specified $\mu$.

\subsection{A free spin without a static polarization}
At $\epsilon=-U/2$ the empty and doubly occupied states both have
energy zero and the two singly occupied states each have energy $-U/2$.
For $T>0$ the grand partition function is $2(1+e^{U/(2T)})$. It
follows that the mean occupation is exactly one at \emph{any}
temperature; it tells us nothing about moment formation. A more
discriminating observable is $4\langle(S^z)^2\rangle$, which equals
the probability of single occupation. The complementary observable,
$\langle(n-1)^2\rangle$, counts empty and doubly occupied sites.
Test the two in a thermal trace over all four charge and spin states:
\begin{wl}
sz2 = pow[number[d[], UP] - number[d[], DO], 2];
q2 = pow[number[d[]] - 1, 2];
Expand[sz2 + q2]
\end{wl}
\begin{wl}
hat = atom /. {\[Epsilon] -> -1, U -> 2};
rows = Table[Prepend[Round[thermalaverage[#, hat, {d}, temp], 1/1000] & /@
    {number[d[]], sz2, q2}, temp], {temp, {0, 1, 4}}];
TableForm[N[rows], TableHeadings -> {None, {"T", "Charge", "Moment", "Charge variance"}}]
\end{wl}
The columns are $T$, $\langle n\rangle$, $4\langle(S^z)^2\rangle$ and
$\langle(n-1)^2\rangle$, for $U=2$ in the energy units used here.
At zero temperature \fn{thermalaverage} averages equally over the two
ground states, so $\langle S^z\rangle=0$ although the spin moment is
maximal. A field term $-h S^z$ lifts that degeneracy. Since $S^z$
commutes with the atomic Hamiltonian, the static susceptibility at
$h=0$ is $\chi=\langle(S^z)^2\rangle/T$ (the magnetic moment multiplying
the field is set to one). At low nonzero $T$ the free-spin contribution
$1/(4T)$ survives; at high $T$ the two spinless states halve its
coefficient. Taking $T\to0$ at fixed $h>0$ selects a pure spin state
with a different magnetization from the zero-field mixture.

\section{Exact solution of the Hubbard dimer}
Connect two sites by \emph{one} bond. The empty and doubly occupied
local states have spin zero, so each can mix with a two-site singlet
but not with a triplet. Build the full Hamiltonian before restricting
the basis. Products of operators are made with \fn{nc}; ordinary
multiplication suffices for a scalar times an operator.
This construction recasts the two-site symmetry problem of Steeb and
Hardy~\cite[Sec.~3.4, problem 84]{steeb} in a basis of total spin
instead of fixed spin projection.
\begin{wl}
sites = {c[1], c[2]};
h0 = U Total[hubbard[sites]];
v = -t hop[c[1], c[2]];
hd = h0 + v;
Expand[{komutator[hd, Total[number /@ sites]],
  komutator[hd, Total[spinplus /@ sites]]}]
\end{wl}
The second zero checks more than conservation of $S^z$: together with
Hermiticity it establishes invariance under spin rotations. At $L=2$,
$Q=N-2$, so the half-filled blocks are those with $Q=0$.
\begin{wl}
bd = qsbasisvc[sites];
sector[s_] := Select[bd, First[#] == {0, s} &];
mat[h_, b_] := First[makematricesbzvc[h, b]][[2]];
hs = mat[hd, sector[0]];
{Dimensions[hs], mat[hd, sector[1]]}
\end{wl}
\begin{wl}
hs // MatrixForm
\end{wl}
In this phase convention the diagonal states, in order, are a doublon
on one site, the spin singlet with one electron on each site, and a
doublon on the other site. Each amplitude $-\sqrt{2}\,t$ is the sum of
two indistinguishable hopping paths out of the normalized singlet.
Rephasing one basis vector flips a row and the corresponding column;
it cannot change an eigenvalue. With $P(x)=\det(H_{S=0}-xI)$:
\begin{wl}
Factor[CharacteristicPolynomial[hs, x]]
\end{wl}
The antisymmetric doublon combination decouples at energy $U$. The
symmetric combination couples to the singly occupied singlet by $-2t$.
The other two singlet energies and the singlet--triplet splitting are
\begin{equation}
 E_\pm=\frac{U\pm\sqrt{U^2+16t^2}}2,\qquad
 \Delta_s=E_T-E_-=-E_- .
 \label{eq:dimer}
\end{equation}
At $U=0$ the singlet fills the bonding orbital with both spins and has
energy $-2|t|$. At $t=0$ the singly occupied singlet and the triplet
are degenerate. Turning on hopping at large $U$ leaves the triplet at
zero but lowers the singlet. This is the simplest example of kinetic
energy selecting antiferromagnetic correlations.

\subsection{What remains of double occupation?}
The derivative $\partial E_-/\partial U=\langle D_1+D_2\rangle$
counts the doublons in the exact ground state. The equality follows by
differentiating the eigenvalue equation for a normalized eigenvector;
terms involving its derivative cancel. It is valid away from a crossing.
\begin{wl}
emin = (U - Sqrt[U^2 + 16 t^2])/2;
dsite = Simplify[D[emin, U]/2];
{dsite, dsite /. {U -> 0, t -> 1}}
\end{wl}
For $U\gg |t|$ the probability per site is $2t^2/U^2$ to leading order.
It is small but cannot be set to zero before evaluating the energy.
In fact $\langle v\rangle=t\,\partial E_-/\partial t$, and at this order
the kinetic gain is twice the interaction cost. Setting $U=4t$ gives
an instructive intermediate case: the exact singlet still has no local
magnetization, even though its local moment
$\langle(n_{i\uparrow}-n_{i\downarrow})^2\rangle=1-2\langle D_i\rangle$
already substantially exceeds the free-electron value $1/2$.

\begin{pitfall}
Occupation vectors refer to the current list of modes. A new call to a
basis generator can change it. Construct the matrix before creating a
basis on different sites, or regenerate the original basis first.
The numerical matrix, once constructed, is independent of that list.
\end{pitfall}

\section{Eliminating virtual charge excitations}
Let $P$ select one electron per site and $\bar P=1-P$. In the half-filled
dimer $P H_0 P=P v P=0$. Every intermediate state connected to $P$ by
one hop has energy $U$, so to second order
\[
 H^{(2)}=-P v\bar P vP/U .
\]
It is negative semidefinite. In a spin triplet there is no accessible
doublon, hence no shift. For a singlet the squared norm of $\bar P v
\ket{s}$ is $4t^2$. Since the singlet projector is
$1/4-\vS_1\cdot\vS_2$, this already fixes the entire spin Hamiltonian.

\begin{wl}
p = nc @@ (projector1 /@ sites);
heff = schriefferwolff[h0, v, p, 2];
snegxoperators[X, "HubbardAtom"];
snegspinoperators[{S, 1/2}];
Expand[xtospin[fermionstox[heff[[3]], c -> X], X -> S]]
\end{wl}

Read the transverse terms with
$S_1^xS_2^x+S_1^yS_2^y=(S_1^+S_2^-+S_1^-S_2^+)/2$.
The result is $J(\vS_1\cdot\vS_2-1/4)$ with $J=4t^2/U$.
\fn{heff[[3]]} is order two: the returned list starts with
order zero. The constant $-J/4$ matters when comparing energies
rather than splittings. A controlled expansion requires $|t|/U\ll1$.
The exact dimer shows both the sign and scale of the first omitted term:
\begin{wl}
Normal[Series[-emin, {t, 0, 4}, Assumptions -> U > 0]] //
  Simplify[#, U > 0] &
\end{wl}
Figure~\ref{fig:dimer-exchange} compares the two truncations. The
fourth-order correction lowers the second-order splitting. At $U=2t$
it cancels the entire leading term, although the exact gap is positive.
The relevant scale is visible already in Eq.~\eqref{eq:dimer}: the
Taylor series of the square root about $t=0$ converges for
$|4t/U|<1$, not throughout the range in the figure.

\begin{figure}[htbp]
\centering
\begin{tikzpicture}
\begin{axis}[width=.95\linewidth,height=53mm,xmin=2,xmax=12,ymin=0,ymax=2.05,
  xlabel={$U/t$},ylabel={$\Delta_s/t$},tick label style={font=\small},
  legend cell align=left,legend style={font=\small,draw=none,fill=none},
  samples=100,domain=2:12]
\addplot[very thick,snegblue] {8/(sqrt(x*x+16)+x)};
\addlegendentry{exact}
\addplot[black,dashed,thick] {4/x};
\addlegendentry{second order}
\addplot[black,densely dotted,thick] {4/x-16/(x*x*x)};
\addlegendentry{fourth order}
\end{axis}
\end{tikzpicture}
\caption{Singlet--triplet gap of a half-filled dimer, $t>0$. All
curves use the same one-bond hopping. At large $U/t$ the remaining
error after the fourth-order correction is of order $(t/U)^5$ in
the plotted units.}\label{fig:dimer-exchange}
\end{figure}

\subsection{Doping: when the first-order term survives}
Remove an electron. The low-energy subspace must now include holes, so
use $P_G=\prod_i(1-D_i)$ rather than the projector onto exactly one
electron per site. A hop onto an empty neighbour need not leave $P_G$.
With $\tilde c_{i\sigma}=c_{i\sigma}(1-n_{i\bar\sigma})$, the
usual nearest-neighbour $t$--$J$ Hamiltonian is
\begin{align}
 H_{tJ}={}&-t\sum_{\langle ij\rangle\sigma}
  (\tilde c^\dagger_{i\sigma}\tilde c_{j\sigma}+\hc)\nonumber\\
 &+J\sum_{\langle ij\rangle}
   (\vS_i\cdot\vS_j-\tfrac14 n_i n_j).
\end{align}
The density factor makes the exchange vanish if either site is empty.
On three or more sites this expression does \emph{not} exhaust
order $t^2/U$. A hop $k\to j$ can create a doublon, followed by a hop
$j\to i$ that removes it with $i\ne k$. If $i$ was initially empty,
the final state is in $P_G$, but the hole has moved from $i$ to $k$.
For real nearest-neighbour hopping the extra term can be written
\[
 -\frac{t^2}{U}\!
 \sum_{j,\;i\ne k\in\mathcal{N}(j),\;\sigma}
 P_G\big(c^\dagger_{i\sigma}n_{j\bar\sigma}c_{k\sigma}
 -c^\dagger_{i\sigma}c^\dagger_{j\bar\sigma}
 c_{j\sigma}c_{k\bar\sigma}\big)P_G .
\]
The neighbours $(i,k)$ are ordered; the reverse path is already in
the sum. The minus sign inside the parentheses distinguishes the two
possible spin transfers. At half filling both transfers to a third site
are blocked by the final projector. Away from half filling, neglecting
them is an additional modelling decision, not a higher-order truncation.

\paragraph{Isolating the three-site matrix element.}
Take an open chain $1$--$2$--$3$ with two electrons in a spin singlet.
There are three states with the electrons on different sites and three
with both electrons on one site. A singlet shared by sites 1 and 2 can
turn into a singlet on 2 and 3 in two hops through a doublon on site 2.
In contrast, exchange on either individual bond cannot change the hole
position. This makes the off-diagonal element of $H^{(2)}$ between
the two bond singlets a diagnostic for the three-site term.
\begin{wl}
line3 = c /@ Range[3];
b3 = Select[qsbasisvc[line3], First[#] == {-1, 0} &];
nd = Diagonal[mat[Total[hubbard[line3]], b3]];
lo = Flatten[Position[nd, 0]]; hi = Flatten[Position[nd, 1]];
v3 = mat[-hop[c[1], c[2]] - hop[c[2], c[3]], b3];
k3 = -v3[[lo, hi]] . v3[[hi, lo]];
k3 // MatrixForm
\end{wl}
The printed matrix is in units of $t^2/U$; the dimensionless entries
are exact. A pair occupying sites 1 and 3 has no virtual doublon one
hop away, hence a zero row. The two adjacent singlets have diagonal
shift $-4$ and mutual coupling of magnitude $2$: only the intermediate
doublon on the shared site connects them. Which row is zero and the
sign of the coupling depend on the chosen basis. The eigenvalues do
not. Removing the three-site term from this matrix gives a
spurious degeneracy:
\begin{wl}
{Sort[Eigenvalues[k3]], Sort[Diagonal[k3]]}
\end{wl}
These numbers diagnose the second-order \emph{operator}, not the full
energy shifts of the doped chain. To obtain the low-energy spectrum
through this order, add the first-order block $t\,v_{3,\mathrm{lo},\mathrm{lo}}$
to $(t^2/U)\,k_3$ and diagonalize. At half filling the first-order
block would be zero.

\subsection{Four sites: a process with no dimer counterpart}
Four successive hops around a square can restore single occupation
while permuting all four spins. On a bipartite plaquette, returning to
single occupation requires an even number of hops, so the first
corrections to superexchange occur at order $t^4/U^3$. They include
both pair exchanges and a cyclic permutation plus its inverse
\cite{macdonald}. Renormalizing $J$ alone cannot describe an entire
plaquette spectrum. A coefficient quoted for a four-spin product is
not directly the coefficient of a permutation: expanding a spin-$1/2$
swap, $P_{ij}=2\vS_i\cdot\vS_j+1/2$, generates lower-body terms too.

The plaquette example compares the lowest sixteen half-filled energies
with the sixteen energies in the projected spin space. From the
checkout run:
\begin{lstlisting}[style=wl]
wolframscript -file examples/schriefferwolff-plaquette.m
\end{lstlisting}
Keeping $t$ fixed and doubling large $U$ reduces the leading
absolute spectral error by about eight at second order ($t^4/U^3$)
and by about thirty-two at fourth order ($t^6/U^5$). Matching the
dimension and the energy origin is essential in this comparison.

\begin{furtherreading}
For the passage from electrons to spins see Auerbach~\cite{auerbach}.
In a junction the virtual states can include superconducting
quasiparticles; their energies replace the single denominator $U$ of
the dimer. An application is the Josephson harmonic expansion of
Bácsi, Iličin and Žitko~\cite{bacsi}.
\end{furtherreading}

\section{Measuring a charge gap on a finite cluster}
Define $E_0(N)$ as the lowest energy among \emph{all} spin sectors
at particle number $N$. At zero temperature that charge sector is stable
over a chemical-potential interval of width
\begin{equation}
 \Delta_c=\mu_+-\mu_-
 = E_0(N+1)+E_0(N-1)-2E_0(N).
 \label{eq:chargegap}
\end{equation}
This second difference eliminates every term linear in $N$. It is
not the energy of the lowest neutral excitation: the dimer, for
instance, admits an arbitrarily soft spin excitation as $t^2/U\to0$,
even as its charge gap grows.

Take four sites on a ring, count each bond once, and measure energies in
units of $t>0$. Here we keep $S^z$, rather than total $S$, and use all
magnetic substates. The blocks are larger but their dimensions now sum
directly to the Fock-space dimension $4^4$.
\begin{wl}
ring = c /@ Range[4];
hr = -Sum[hop[c[i], c[Mod[i, 4] + 1]], {i, 4}] +
  U Total[hubbard[ring]];
br = qszbasisvc[ring];
blocks = makematricesbzvc[hr, br];
{Length[blocks], Total[Length /@ blocks[[All, 2]]]}
\end{wl}
\begin{wl}
e0[q_, u_] := Min[Flatten[Eigenvalues[N[#[[2]] /. U -> u]] & /@
  Select[blocks, #[[1, 1]] == q &]]];
gap[u_] := e0[1, u] + e0[-1, u] - 2 e0[0, u];
rows = Table[{u, N[Round[gap[u], 1/1000]]}, {u, {0, 2, 4, 8, 16}}];
TableForm[rows, TableHeadings -> {None, {"U/t", "Charge gap / t"}}]
\end{wl}
Each row gives $(U/t,\Delta_c/t)$ rounded to three decimal places.
At $U=0$ the allowed momenta are $0,\pm\pi/2,\pi$ (lattice spacing
one). The two zero-energy orbitals make this an open shell, so adding
and removing an electron have identical one-particle thresholds.
The same calculation on a closed-shell cluster can yield a nonzero
gap in a metal. Inferring a bulk gap requires a sequence of sizes
and control of such shell effects.

For comparison, at $U=0$ alternate the hopping amplitudes
$t_1,t_2>0$ in a chain with two sites per cell. The Bloch energies are
$\pm|t_1+t_2 e^{ik}|$, where $k$ is in inverse cell lengths. The
half-filled band gap $2|t_1-t_2|$ persists in the thermodynamic limit,
without an interaction. In the large-$U$ uniform Hubbard chain charge
motion instead costs a doublon energy even in the absence of a
one-particle band gap. The charge-gap measurement alone cannot tell
these mechanisms apart; compare the noninteracting spectrum, double
occupancy and spin response.

\subsection{A shell-effect check without many-body diagonalization}
The noninteracting uniform ring is a useful calibration. Sort its
$L$ one-particle levels \emph{per spin}; at half filling each spin
occupies $L/2$ levels. The finite-size addition gap is the difference
of two consecutive levels at the Fermi surface.
\begin{wl}
shell[l_] := With[{band = SortBy[Table[-2 Cos[2 Pi j/l], {j, 0, l-1}], N]},
  band[[l/2 + 1]] - band[[l/2]]];
TableForm[Table[{l, N[shell[l]]}, {l, {4, 6, 8, 10}}],
  TableHeadings -> {None, {"Sites", "Charge gap / t"}}]
\end{wl}
For $L=4p$, both Fermi points lie on the grid and the gap is zero.
For $L=4p+2$, the grid misses each Fermi point by $\pi/L$, so
$\Delta_c/t=4\sin(\pi/L)$: a finite gap at \emph{every} size on this
subsequence tends to zero as $1/L$. In an interacting calculation one
cannot simply subtract this formula; the velocities and the many-body
levels have changed. But treating the two subsequences separately
prevents an alternation of boundary-condition effects from masquerading
as an interaction threshold.

\section{The atomic propagator and Hubbard-I}
Consider the retarded one-electron Green's function
$G_\sigma(z)=\langle\!\langle d_\sigma;d^\dagger_\sigma
\rangle\!\rangle_z$, where $z=\omega+i0^+$. Its spectral density
$A_\sigma(\omega)=-\operatorname{Im}G_\sigma(\omega+i0^+)/\pi$
has unit integral \emph{per spin}. Use the same grand Hamiltonian
as in the atomic spectrum. Since $n_{\bar\sigma}$ commutes with it,
the two operators $d_\sigma(1-n_{\bar\sigma})$ and
$d_\sigma n_{\bar\sigma}$ evolve at fixed frequencies $\epsilon$ and
$\epsilon+U$, respectively. The equation-of-motion hierarchy closes:
\begin{align*}
 (z-\epsilon)G_\sigma &= 1+U F_\sigma,\\
 (z-\epsilon-U)F_\sigma &= \langle n_{\bar\sigma}\rangle ,
\end{align*}
where $F_\sigma=\langle\!\langle d_\sigma n_{\bar\sigma};
d^\dagger_\sigma\rangle\!\rangle_z$. These are operator identities;
no factorization of an average was used.
\begin{wl}
gatom = greensfunction[d[AN, UP], d[CR, UP], atom, z];
g = gatom /. avg[number[d[], DO]] -> w;
Apart[g, z]
\end{wl}
Writing $w=\langle n_{\downarrow}\rangle$, the poles and their weights
are transparent in
\begin{equation}
 G_\uparrow(z)=\frac{1-w}{z-\epsilon}
              +\frac{w}{z-\epsilon-U}.
 \label{eq:atomicG}
\end{equation}
The upper pole exists only when a down electron is available on the
site. With $\epsilon=-U/2$ and a spin-symmetric ensemble, $w=1/2$;
the spectral peaks then have equal weight and are separated by $U$.
A pure member of the atomic doublet has the same energy but different
spin-resolved weights. Substituting $w=1/2$ is a specification of the
state, not a step in the equation-of-motion calculation.
\begin{wl}
ph = Factor[g /. {\[Epsilon] -> -U/2, w -> 1/2}];
{ph, Limit[z ph, z -> Infinity]}
\end{wl}

\paragraph{Checking both residues against a trace.}
The half-filled atom is a poor test of thermal weights: symmetry
forces them both to $1/2$ regardless of temperature. Instead take
$\epsilon=-1/4$, $U=2$, and $T=1/2$. Removing or adding one spin-up
electron still gives only the two transition frequencies $\epsilon$
and $\epsilon+U$, but their residues differ. The Lehmann representation
sums matrix elements between exact many-body eigenstates. Each row
below is $(E_n-E_m,w_{mn})$ after collecting equal frequencies:
\begin{wl}
asym = atom /. {\[Epsilon] -> -1/4, U -> 2};
poles = lehmann[d[AN, UP], d[CR, UP], asym, {d}, 1/2];
TableForm[N[Round[poles, 1/1000]], TableHeadings -> {None, {"Frequency", "Residue"}}]
\end{wl}
For a retarded anticommutator the factor in the residue is the
\emph{sum} of the two Boltzmann probabilities, not their difference.
The atomic partition function is
$Z=1+2e^{-\epsilon/T}+e^{-(2\epsilon+U)/T}$.
The upper peak has weight $w=(e^{-\epsilon/T}+e^{-(2\epsilon+U)/T})/Z$.
These facts give two independent checks: the total weight is one,
and the upper weight equals the mean down-spin occupation. Compare the
full complex propagators at a frequency off the real axis, where a
pole neither diverges nor needs artificial broadening:
\begin{wl}
wth = thermalaverage[number[d[], DO], asym, {d}, 1/2];
z0 = 1/3 + I/4;
gtrace = Total[poles[[All, 2]]/(z0 - poles[[All, 1]])];
{Abs[1 - Total[poles[[All, 2]]]] < 10^-12,
 Abs[gtrace - (g /. {z -> z0, \[Epsilon] -> -1/4, U -> 2, w -> wth})] < 10^-12}
\end{wl}

\subsection{What changes when the atom is put back in the lattice?}
Define the local self-energy by
$G_\mathrm{at}^{-1}(z)=z-\epsilon-\Sigma_\mathrm{at}(z)$. Inverting
Eq.~\eqref{eq:atomicG} gives
\[
 \Sigma_\mathrm{at}(z)=Uw+
 \frac{U^2 w(1-w)}{z-\epsilon-U(1-w)} .
\]
The static term is the Hartree shift. The pole in the second term is
what opens a correlation gap in this approximation. Hubbard-I inserts
this atomic self-energy into the lattice propagator, or equivalently
uses $G_k(z)=[G_\mathrm{at}^{-1}(z)-\varepsilon_k]^{-1}$, with
$\varepsilon_k$ the dispersion of the hopping term alone.
\begin{wl}
gk = Together[ph/(1 - ek ph)];
Factor[Denominator[gk]]
\end{wl}
At particle--hole symmetry its poles are
$z_\pm=(\varepsilon_k\pm\sqrt{\varepsilon_k^2+U^2})/2$.
For $U>0$ they lie on opposite sides of zero. For $U=0$ the apparent
additional pole cancels against a zero in the numerator, leaving
$1/(z-\varepsilon_k)$. Checking residues as well as pole locations
prevents a spurious flat band. The construction is exact in the atomic
limit and recovers free propagation at $U=0$. It omits intersite
spin correlations and has no coherent Fermi-liquid resonance. Its
half-filled gap at arbitrarily small $U>0$ does not establish a Mott
transition for a general lattice.

\section{Can a static field account for the singlet energy?}
In a collinear Hartree state, write $n_{i\sigma}=\bar n_{i\sigma}
+\delta n_{i\sigma}$ and discard
$\delta n_{i\uparrow}\delta n_{i\downarrow}$. The remaining interaction is
\begin{equation}
 U\sum_i(\bar n_{i\downarrow} n_{i\uparrow}+
          \bar n_{i\uparrow} n_{i\downarrow}-
          \bar n_{i\uparrow}\bar n_{i\downarrow}).
\end{equation}
The last term does not change the orbitals, but omitting it counts
the interaction energy twice. In a translationally invariant metal,
a small uniform polarization $m=\bar n_\uparrow-\bar n_\downarrow$
splits the two spin bands by $Um$. For a smooth one-spin density of
states $\rho(E_F)$ per site, the induced polarization at zero
temperature is $\rho(E_F)Um$ to first order. Thus the unpolarized
solution is unstable to \emph{uniform} magnetization for
$U\rho(E_F)>1$ (Stoner criterion). It need not be the first
instability: the susceptibility at nonzero wave vector can be larger.

On the dimer, try staggered polarization instead:
$m_1=-m_2=m$, with one electron per site on average. Apart from the
common Hartree shift, each spin sees diagonal potentials $\pm Um/2$
and off-diagonal hopping $-t$. Filling the lower eigenstate for each
spin gives
\[
 m=\frac{Um/2}{\sqrt{t^2+(Um/2)^2}} .
\]
Besides $m=0$, for $U>2|t|$ there are two solutions related by spin
reversal with $m=\pm\sqrt{1-(2t/U)^2}$. Below we seed one branch;
the solver adjusts the chemical potential for two particles.
\begin{wl}
hf = meanfield[hd /. {t -> 1, U -> 4}, "Hartree"];
seed = {avg[number[c[1], UP]] -> 0.9, avg[number[c[2], DO]] -> 0.9};
sol = meanfieldsolve[hf, sites, seed, "ParticleNumber" -> 2];
m1 = avg[number[c[1], UP]] - avg[number[c[1], DO]];
{sol["Converged"], Round[sol["Energy"], 1/1000],
  Round[m1 /. sol["Averages"], 1/1000] // N}
\end{wl}
At these parameters the unpolarized determinant has energy
$-2t+U/2=0$, the staggered determinant has energy $-t/2$, and the
exact energy is $(2-2\sqrt{2})t$. Improving the variational energy by
allowing a staggered field has introduced a false broken symmetry on a
finite dimer. In particular $m_1$ is not the local fluctuating moment
calculated in the singlet. At large $U$, Hartree theory yields
$-2t^2/U$ for this determinant; the exact singlet gains twice as much
through transverse spin exchange. Convergence of the iteration is not
a test of either the symmetry or the accuracy of the state.

\section{Local constraints: Hubbard operators and auxiliary particles}
Operator projection is easiest to understand on one site. For
$a,b\in\{0,\uparrow,\downarrow,2\}$ let $X^{ab}=\ket{a}\bra{b}$.
Then $X^{ab}X^{cd}=\delta_{bc}X^{ad}$ and $\sum_a X^{aa}=1$.
The creation operator $d^\dagger_\uparrow$ is the sum of the transitions
$X^{\uparrow0}+X^{2\downarrow}$. For down spin the second transition
has a minus sign because $\ket{2}$ was defined with the up creator
on the left. Work with the corresponding fermionic polynomials to
test the projected algebra:
\begin{wl}
xi = xoperator[d[], "0", "UP"];
Expand[antikomutator[xi, conj[xi]]]
\end{wl}
The anticommutator is $X^{00}+X^{\uparrow\uparrow}
=1-n_\downarrow$, not the identity. The missing weight is precisely
the transition at the upper Hubbard pole. At distinct sites, odd
transitions anticommute: ordinary commuting site matrices would lose
the fermion sign.
\begin{wl}
pg = 1 - hubbard[d[]];
Expand[{nc[pg, pg] - pg, nc[pg, d[AN, UP], pg] - xi}]
\end{wl}

One can instead keep canonical auxiliary fields at the price of a
constraint. In the no-doublon space introduce a boson $b$ for
$\ket0$ and fermions $f_\sigma$ for $\ket\sigma$. Identify
$X^{0\sigma}=b^\dagger f_\sigma$ and impose
\[
 b^\dagger b+\sum_\sigma f^\dagger_\sigma f_\sigma=1
 \quad\text{on every site}.
\]
The three physical states are precisely the states with one auxiliary
particle. The substitutions $b\to e^{i\theta}b$ and
$f_\sigma\to e^{i\theta}f_\sigma$ leave $X^{0\sigma}$ invariant:
the common phase is redundant. Discarding the constraint introduces,
among others, an auxiliary vacuum and a two-spinon state with no
counterparts in the projected atom. Enforcing it only on average is
a mean-field approximation. The idempotent $P_G$ and the Hubbard
operators above implement the projection exactly and offer local
tests for any auxiliary-particle calculation.

\subsection{Partial projection as a variational wave function}
What is lost if one discards doublons altogether? Return to the dimer,
with $t>0$ and exactly two particles. The free bonding determinant is
an equal-amplitude superposition of the singly occupied singlet
$\ket{s}$ and the symmetric doublon $\ket{D_+}$, in phases chosen
so that $\langle s|H|D_+\rangle=-2t$. Multiply each configuration
with a doublon by $g$, $0\le g\le1$. This is a small-system version
of a Gutzwiller variational state:
\[
 \ket{\psi(g)}=\frac{\ket{s}+g\ket{D_+}}{\sqrt{1+g^2}},
 \qquad E(g)=\frac{-4tg+Ug^2}{1+g^2}.
\]
There is no approximation in evaluating this expectation value.
More unusually, on the dimer this one-parameter family contains the
exact ground state. Stationarity gives $2tg^2+Ug-2t=0$, whose
positive root can be written without subtracting nearly equal numbers:
\begin{wl}
gopt = 4 t/(U + Sqrt[U^2 + 16 t^2]);
energy[y_] := (-4 t y + U y^2)/(1 + y^2);
FullSimplify[{energy[gopt] - emin,
  gopt^2/(2 (1 + gopt^2)) - dsite}, U > 0 && t > 0]
\end{wl}
The second check is the probability of finding a doublon on one
specified site. At $t/U\ll1$ the optimum has $g\simeq2t/U$, whereas
strict projection $g=0$ gives $E=0$ when evaluated with the original
Hamiltonian. A nonzero admixture of \emph{amplitude} $t/U$ produces
exchange energy of order $t^2/U$. The projected spin Hamiltonian
recovers that energy by transforming the Hamiltonian, rather than
by retaining the doublon in the state. For larger clusters inequivalent
hole--doublon separations generally require different amplitudes; one
scalar $g$ is no longer enough to make the trial state exact.

\section{A three-orbital atom: Hund's rule from the spectrum}
Distinguish a single-particle level index from a site index: now
$a_{m\sigma}$ removes an electron in orbital $m=1,2,3$ of \emph{one}
atom. Set all three orbital energies to zero. For a rotationally
invariant Kanamori interaction the repulsion in distinct orbitals is
$U'=U-2J$ for opposite spin and $U'-J$ for equal spin. Density terms
alone would give different energies to the three components of a spin
triplet. Two additional processes restore the multiplet structure:
\begin{align}
 H_K={}&U\sum_m D_m+
 (U-2J)\sum_{m\ne n} n_{m\uparrow}n_{n\downarrow}
 +(U-3J)\sum_{m<n,\sigma} n_{m\sigma}n_{n\sigma}\nonumber\\
 &-J\sum_{m\ne n} S_m^+ S_n^-
 +J\sum_{m\ne n} B_m^\dagger B_n,\qquad
 B_m=a_{m\downarrow}a_{m\uparrow}.
 \label{eq:kanamori}
\end{align}
The ordered last sums already contain their Hermitian conjugates. The
pair convention matters: reversing only the two annihilators would
reverse the pair-hopping amplitude. To make all factors explicit,
construct the terms for each unordered orbital pair:
\begin{wl}
snegfermionoperators[a];
snegrealconstants[J];
orb = a /@ Range[3];
n[m_, s_] := number[a[m], s];
pair[m_] := nc[a[AN, m, DO], a[AN, m, UP]];
hkpair[m_, k_] := (U - 2 J) (nc[n[m, UP], n[k, DO]] +
    nc[n[m, DO], n[k, UP]]) +
  (U - 3 J) Total[nc[n[m, #], n[k, #]] & /@ {UP, DO}] -
  J (nc[spinplus[a[m]], spinminus[a[k]]] +
     nc[spinminus[a[m]], spinplus[a[k]]]) +
  J (nc[conj[pair[m]], pair[k]] + nc[conj[pair[k]], pair[m]]);
hk = U Total[hubbard[orb]] + Total[hkpair @@@ Subsets[Range[3], {2}]];
Expand[{hk - conj[hk], komutator[hk, Total[spinplus /@ orb]]}]
\end{wl}

Put two electrons in six spin-orbitals. Here $Q=N-3=-1$. The
following rows give $S$, the energy (in symbolic $U,J$), and the
number of \emph{multiplets} at that energy. We keep the multiplicities;
distinct eigenvalues alone cannot diagnose a missing state.
\begin{wl}
bk = qsbasisvc[orb];
k2 = makematricesbzvc[hk, Select[bk, #[[1, 1]] == -1 &]];
levels = Flatten[Table[{b[[1, 2]], #[[1]], #[[2]]} & /@
    Tally[Expand /@ Eigenvalues[b[[2]]]], {b, k2}], 1];
TableForm[Sort[levels], TableHeadings -> {None, {"Spin", "Energy", "Multiplets"}}]
\end{wl}
For $J>0$ the lowest energy is $U-3J$. The two electrons occupy
different orbitals with parallel spins; three choices of orbital pair
give three triplets. Among singlets, each distinct-orbital pair gives
one state at $U-J$. The three intra-orbital pairs form a matrix with
diagonal $U$ and off-diagonal $J$. Its symmetric vector has energy
$U+2J$, and the two orthogonal vectors have energy $U-J$. Thus the
dimensions are $9+5+1=\binom62$. Dropping pair hopping would leave
three degenerate states at $U$ instead of this $1+2$ splitting.
\begin{wl}
Total[(2 #[[1]] + 1) #[[3]] & /@ levels] == Binomial[6, 2]
\end{wl}
\begin{furtherreading}
The atomic splittings determine which spin and orbital degrees of
freedom survive when the shell couples to a bath. For the three-orbital
impurity with this interaction see Horvat, Žitko and Mravlje~\cite{horvat}.
The \href{https://auger.ijs.si/nano/biblio/biblio.html}{authors' publication list}
provides access to the paper.
\end{furtherreading}

\section*{Exercises}
\addcontentsline{toc}{section}{Exercises}
\begin{exercise}[Counting charge defects]\label{ex:defects}
At $t=0$ count the half-filled ground states and the states exactly
$U$ higher on $L\ge2$ sites. How many states must a charge-only basis
contain at each energy for $L=4$?
\end{exercise}

\begin{exercise}[Detuned sites]\label{ex:detuned}
Add $\delta(n_2-n_1)/2$ to the half-filled dimer. Derive the exchange
$J(\delta)$ to second order in $t$ and state the condition on
$U,\delta,t$ for this spin-only description. Compare with a singlet
eigenvalue, taking $U=10$, $\delta=6$ and $t=1/10$ in common energy units.
\end{exercise}

\begin{exercise}[Addition versus spin gap]\label{ex:gaps}
At half filling find the dimer charge gap for all $U\ge0$, $t>0$.
Why does it remain positive at $U=0$? Compare it with the spin gap
at $U=0$ and at $U/t\gg1$.
\end{exercise}

\begin{exercise}[A spectral sum rule]\label{ex:moments}
Derive the first three high-frequency coefficients of the atomic
propagator at general filling. Which of them distinguish the two
Hubbard poles from a single Hartree pole? Verify with nested
commutators without extracting the poles a second time.
\end{exercise}

\chapter{Superconductivity}
\setwlchapter{superconductivity}
\section{From a Cooper pair to the gap equation}
Above an inert Fermi sea two electrons with opposite momenta can
scatter repeatedly within a shell of width $\omega_D$. For a constant
attraction $-g$ (with the volume absorbed in the normalization of the
sum), their binding energy $E_b>0$ satisfies
$1=g\sum_{0<\xi_k<\omega_D}(2\xi_k+E_b)^{-1}$. A nonzero one-spin
density of states at the Fermi surface makes this sum logarithmic as
$E_b\to0$. The instability is at arbitrarily weak attraction; its
binding energy is not a power series in $g$.

In a condensate one instead decouples the reduced interaction with
$\Delta=g\sum_k\langle c_{-k\downarrow}c_{k\uparrow}\rangle$.
For real positive $\Delta$ the pair-creation term has coefficient
$-\Delta$. Diagonalizing the particle--hole block gives
\[
 E_k=\sqrt{\xi_k^2+\Delta^2},\qquad
 |v_k|^2=(1-\xi_k/E_k)/2,\qquad |u_k|^2=1-|v_k|^2 .
\]
The gap equation contains the thermal factor $1-2f(E_k)$.
With the sum now covering both signs of $\xi_k$, the weak-coupling
gap exponent differs by a factor of two from the isolated Cooper
pair exponent. A band cutoff, the number of spin species in the
density of states, and the normalization of $g$ must agree before
comparing the two.
\casestudy{gap}{Self-consistency from the anomalous average}

\section{Nambu doubling and bound states}
For a single spinful orbital there are two physical positive-energy
quasiparticles, not four. The negative eigenvalues of the BdG matrix
are their particle--hole partners. A bound state can cross zero as a
function of phase or exchange coupling. Its occupation and the
many-body ground-state parity then matter even if a plot only shows
positive energies. In an interacting dot at large lead gap the
induced pairing mixes $|0\rangle$ and $|2\rangle$; the two singly
occupied states remain a spin doublet. Comparing the lower even
eigenvalue to the odd one locates a crossing without broadening.
\casestudy{andreev}{An interacting level with induced pairing}
For an ideal short single-channel junction with normal transmission
$\tau$ and equal lead gaps, $E_A=\Delta\sqrt{1-\tau\sin^2(\phi/2)}$.
At zero temperature the even ground branch contributes
$I=(2e/\hbar)\partial_\phi(-E_A)$; counting the two Nambu partners
as independent filled levels would double this current. A local
magnetic exchange can also bind a quasiparticle below the gap.
The induced-pairing dot is a useful limiting model, but it does not
retain the continuum at $|\omega|=\Delta$ or its gap-edge singularity.
\subsection{Reading the current from the parity branches}
For a dot in the infinite-gap limit, amplitudes for transferring a
pair through the two contacts add coherently. In a gauge where the
left contact is real, $\Gamma_\phi=\Gamma_L+\Gamma_R e^{i\phi}$.
At particle--hole symmetry the lowest even energy is $-|\Gamma_\phi|$;
the odd energy is $-U/2$ and has no phase dependence in this limit.
The equilibrium current switches branches at a ground-state crossing.
Fixing the parity experimentally instead follows one branch through
the crossing. Here $\Gamma_L=1$, $\Gamma_R=3/5$ and $U=2$.
\begin{wl}
ee[p_] := -Sqrt[1+(3/5)^2+2 (3/5) Cos[p]];
eo = -1;
TableForm[Table[{p, N[ee[p]], eo,
  If[ee[p] < eo, N[(3/5) Sin[p]/(-ee[p])], 0]}, {p,{0,Pi/3,2 Pi/3,Pi}}],
 TableHeadings -> {None,{"Phase","Even","Odd","Current / I0"}}]
\end{wl}
The current scale is $I_0=2e\Gamma_L/\hbar$. At a crossing the
ground energy need not have a derivative; take one-sided derivatives
rather than averaging them. A nonzero odd-branch current at finite
lead gap is not captured by the static model.
\begin{furtherreading}
For the role of the superconducting continuum in variational states
and in Andreev spin qubits see Refs.~\cite{ilicin-variational,ilicin-qubit}.
The residual quasiparticle in a strongly coupled junction is discussed
in Ref.~\cite{pavesic-residual}.
\end{furtherreading}

\section{Odd-parity pairing and end modes}
For a uniform spinless chain the particle--hole diagonal term is
$-\mu-2t\cos k$ and nearest-neighbour odd pairing has magnitude
$2|\Delta\sin k|$. A bulk gap closes at $k=0$ or $\pi$ when
$|\mu|=2|t|$ for nonzero $\Delta$. An open chain in the intervening
phase has end Majoranas; finite chains generally have an exponentially
small splitting, not an exactly protected zero eigenvalue at every
length.
\casestudy{kitaev}{The bulk closing at the topological transition}

\section{Pair pseudospin on a bipartite lattice}
The operator $\eta^+=\sum_j(-1)^j c^\dagger_{j\uparrow}c^\dagger_{j\downarrow}$
creates an on-site pair with staggered phase. It commutes with nearest
neighbour bipartite hopping: the two ends of a bond contribute with
opposite signs. The local charge term gives
$[H-\mu N,\eta^+]=(U-2\mu)\eta^+$, so at half filling it generates a
charge SU(2) alongside ordinary spin SU(2). This exact symmetry of the
Hubbard Hamiltonian is distinct from introducing a BCS anomalous
average.
\begin{exercise}[The atomic parity crossing]\label{ex:andreev}
For one level with repulsion $U$ and induced pairing amplitude
$\Gamma>0$, put $\epsilon=-U/2$. Find the two even energies and the
odd energy. At what $U/\Gamma$ is the lowest state a doublet?
\end{exercise}

\chapter{Magnetism}
\setwlchapter{magnetism}
\section{Exchange on a cluster}
A positive $J$ in $J\vS_i\cdot\vS_j$ favours a singlet. Superexchange
from a half-filled one-band model has this sign, whereas a direct
Hund interaction favours alignment. In double exchange a mobile
carrier aligned with local moments instead gains kinetic energy
when neighbouring moments are parallel. These mechanisms need not
produce the same dependence on filling or on bond angle.

For three spins $1/2$ on an equilateral triangle all pair bonds are
present, so the energy depends only on the total spin:
$H=J[\mathbf S_\mathrm{tot}^2-9/4]/2$. Its two low-spin doublets
have the same energy even though no state is a singlet on all three
bonds at once. On a square ring not all pairs are neighbours;
the Casimir alone no longer fixes each energy.
\casestudy{frustration}{Resolving the four-site ring by magnetization}

\section{Spin waves and their regime of validity}
Expanding around a ferromagnet polarized along $+z$ gives
$S^z=s-b^\dagger b$, $S^+\simeq\sqrt{2s}\,b$. In
$-J\sum_i\vS_i\cdot\vS_{i+1}$ the quadratic magnon energy is
$2Js(1-\cos k)$; the absent gap follows from the global spin-rotation
symmetry. For a Néel reference state rotate one sublattice first.
Opposite classical spins then become parallel local axes, and the
transverse exchange generates \emph{pairs} of bosons. Diagonalizing
the bosonic quadratic form requires the commutation metric with
opposite signs on particles and holes; a unitary diagonalization of
the Hermitian energy matrix does not give the magnon frequencies.
\casestudy{magnons}{A staggered field as a gap regulator}
The quadratic truncation assumes a small mean occupation relative to
$2s$. A gapless antiferromagnetic chain of spin $1/2$ is particularly
unfavourable for a broken-symmetry expansion; finite positive field
in the example is part of the Hamiltonian, not just a numerical tolerance.

\section{One dimension: strings and boundaries}
The Jordan--Wigner string counts the occupied sites to the left of
each spin flip. Adjacent strings cancel in the bulk. Long-range flips
retain intervening density factors, and closing a chain into a ring
introduces a boundary sign that depends on total fermion parity.
The open XXZ chain is therefore a useful spectrum-level test before
handling periodic boundary conditions.
\casestudy{xxz}{The interacting spinless chain}

\section{Itinerant and single-ion magnetism}
In a uniform collinear Hartree state, an up electron feels $U n_\downarrow$.
Linearizing the spin-density difference gives $m=U\rho(E_F)m$ at zero
temperature for a smooth one-spin density of states. A finite cluster
needs a thermal derivative rather than a continuum density at a
single energy. Passing this uniform instability says nothing about
the preferred ordering wave vector.
\casestudy{stoner}{A fixed-density susceptibility test}
For a magnetic ion with $s\ge1$, $D(S^z)^2$ can split a multiplet even
in zero field. A term $E[(S^x)^2-(S^y)^2]$ mixes states whose $m_z$
differs by two. Stevens operators preserve the connection between
crystal symmetry and allowed tensor components. For a bond lacking
inversion symmetry the antisymmetric exchange
$\mathbf D_{ij}\cdot(\vS_i\times\vS_j)$ is allowed;
$\mathbf D_{ji}=-\mathbf D_{ij}$ is required by reversing the bond.
It is a two-site coupling, not a single-ion anisotropy.
\casestudy{anisotropy}{Zero-field splitting of a spin one}
\begin{exercise}[A frustrated triangle]\label{ex:triangle}
Find the energies and degeneracies of the Heisenberg triangle at
$J>0$. Reconstruct its zero-field partition function; the entropy
at $T\to0$ counts the full manifold, not just highest-weight states.
\end{exercise}

\chapter{Quantum impurities and nanostructures}
\setwlchapter{impurities}
\section{A local level embedded in a band}
For one spinful impurity the Anderson Hamiltonian reads
\[
 H=\sum_{k\sigma}\epsilon_k n_{k\sigma}+\epsilon_d n_d+U D_d+
 \sum_{k\sigma}(V_k c^\dagger_{k\sigma}d_\sigma+\hc).
\]
At $V_k=0$ the local moment exists for $\epsilon_d<0<\epsilon_d+U$.
A bath adds infinitely many low-energy configurations; it cannot in
general be represented by a constant shift of the impurity level.
Already at $U=0$ the retarded hybridization is
$\Delta(z)=\sum_k|V_k|^2/(z-\epsilon_k)$, whose imaginary part
describes the widths of levels only in the continuum limit.
\casestudy{leads}{Hybridization to a local lattice orbital}
\casestudy{resonance}{Frequency-dependent hybridization with one bath level}

\section{From valence fluctuations to Kondo exchange}
For initial and final bath states close to the Fermi level, eliminate
empty and doubly occupied impurity states in second order. For real
constant $V$ the exchange multiplying $\vS_d\cdot\mathbf s(0)$ is
$J=2V^2[1/(-\epsilon_d)+1/(\epsilon_d+U)]$. Both virtual paths favour
the impurity--bath singlet. Their \emph{difference} controls potential
scattering, which vanishes at $\epsilon_d=-U/2$. The spin density here
uses Pauli matrices divided by two; changing that definition changes
the quoted $J$. This elimination is controlled when both charge
excitation energies are large on the hybridization scale. It does not
directly compute the nonperturbative Kondo scale.
\casestudy{anderson}{Checking exchange against exact diagonalization}

\section{Wilson chain and NRG Ljubljana}
A uniform finite bath has a lowest level spacing that competes with
the screening scale. NRG instead divides the band into logarithmic
intervals and transforms the retained combinations into a chain.
Only the first chain orbital couples to the impurity. Far along a
metallic Wilson chain the hopping typically decays as
$\Lambda^{-n/2}$ for logarithmic discretization parameter $\Lambda>1$.
Successive diagonalizations can then discard high-energy states relative
to the next hopping. Check a discretization and truncation sequence;
a two-site singlet in an untruncated test is a normalization check,
not evidence by itself of continuum screening.

In the public \href{https://github.com/rokzitko/nrgljubljana}{NRG Ljubljana}
tree, \path{test/c++/test42_model_1orb/model.m} and its companion
\texttt{param} give a complete one-channel initializer example.
\fn{def1ch[1]} sets up the impurity and bath end; the local term uses
\fn{number[d[]]} and \fn{hubbard[d[]]}, the contact uses
\fn{hop[f[0],d[]]}, and the host term is included once. The initializer
derives the contact normalization from the discretized bath; inserting
the continuum width $\Gamma$ as the hopping would have the wrong
normalization. The setting \texttt{symtype=QS} assumes charge U(1) and
spin SU(2) invariance; choose different
quantum numbers if the model has pairing, spin anisotropy or a magnetic
field. Specify measured observables too: an energy spectrum contains
no information about matrix elements of a local spin flip.

\section{Charge stability and superconducting islands}
In a double dot the electrostatic term $U_{12}n_1 n_2$ competes
with the on-site $U$; near the $(1,1)$ configuration the virtual
charge-transfer cost is $U-U_{12}$ for equivalent sites. The same
second-order hopping that mixes $(2,0)$ and $(0,2)$ with a singlet
therefore controls the singlet--triplet splitting. Close to a charge
degeneracy both charge configurations must be retained.
For multi-orbital dots a high-spin atomic multiplet can replace the
spin-$1/2$ moment (see Chapter~10 and Ref.~\cite{horvat}).

A superconducting island brings both pairing and charging energy
$E_C(N-N_g)^2$. A static pair field can model proximity to a bulk
reservoir but is not a number-conserving description of a floating
island. Compare even and odd \emph{total} parity sectors, include
the energy of a residual quasiparticle, and distinguish a spin's
local response from its total-spin quantum number. The latter
distinction enters the impurity Knight shift~\cite{pavesic-knight};
crystal fields provide another way of altering the impurity response
without changing its charge~\cite{zitko-anisotropic}.
\begin{exercise}[Where does the pole weight go?]\label{ex:resonance}
Hybridize a noninteracting level at zero energy with one bath level
also at zero energy by a real $V>0$. Find both poles of the impurity
propagator, their residues per spin, and the limit as $V\to0$.
\end{exercise}

\chapter{Fermi liquids, Green's functions and response}
\setwlchapter{response}
\section{Resolvent or thermal trace?}
There are two ways to compute a retarded propagator on a finite
cluster. Commuting $A$ with $H$ generates a space of operators.
If it closes with $[O_i,H]=\sum_j M_{ij}O_j$, solving a linear
system at $z=\omega+i\eta$, $\eta>0$, sums all nested commutators.
Alternatively, a Lehmann sum uses the many-body energies and
transition matrix elements. The first route can remain symbolic;
the second is a numerical reference that includes correlations
without a factorization.
\casestudy{propagation}{How large an operator space does a chain require?}
\casestudy{lehmann}{Positive diagonal weight versus signed cross weight}
\section{High-frequency checks}
For an electron the anticommutator fixes the $1/z$ coefficient of
the diagonal propagator to one. The next moment probes the mean
addition energy, and a higher moment can distinguish a single pole
from a pair of poles with the same centroid. For a conserved density
the ordinary commutator with itself is zero: a nonzero thermal
fluctuation need not produce a nonzero dynamical retarded response.
\casestudy{moments}{An independent sum rule without finding the poles}

\section{A static limit is a prescription}
Define $\chi^R_{AB}(t)=-i\theta(t)\langle[A(t),B]\rangle$.
With a perturbation $+h B$, $\delta\langle A\rangle=\chi^R_{AB} h$
in this convention; the common Zeeman coupling is instead $-h S^z$.
For the density of a free band the retarded bubble (spin sum included)
is
\[
 \chi^0(q,z)=\frac{1}{N_s}\sum_{k\sigma}
 \frac{f(\epsilon_k)-f(\epsilon_{k+q})}{z+\epsilon_k-\epsilon_{k+q}} .
\]
At $q=0$ and finite $z$ the numerator vanishes. Taking the static
long-wavelength limit instead gives the negative thermodynamic density
of states for a perturbation raising the level energy. In the random
phase approximation a repulsive density interaction $v_q$ induces an
additional potential $v_q\delta n$, giving
$\chi=\chi^0/(1-v_q\chi^0)$. This argument fixes the denominator sign
without importing a convention for a positively defined polarization.
A current response to a vector potential also has a diamagnetic term:
the current itself depends on the field. Omitting it can produce a
spurious response to a pure gauge~\cite{vucicevic}.
\subsection{A two-level susceptibility by two independent routes}
For one electron of fixed spin on a dimer with hopping $-t$, the
bonding state has energy $-t$. The imbalance $D=n_1-n_2$ connects
it to the antibonding state at $+t$ with unit matrix element. Thus
$\chi^R_{DD}(z)=1/(z-2t)-1/(z+2t)$. Unlike total charge, this density
does not commute with the hopping. Check the static value at $t=1$
by applying the field directly. In the site basis $+hD$ changes the
diagonal energies in opposite directions:
\begin{wl}
dop = DiagonalMatrix[{1,-1}];
hd[h_] := {{h,-1},{-1,-h}};
{vals,vecs} = Eigensystem[hd[0]];
j = First[Ordering[vals,1]]; l = 3-j;
delta = vals[[l]]-vals[[j]];
wgt = Abs[Conjugate[Normalize[vecs[[j]]]].dop.Normalize[vecs[[l]]]]^2;
TableForm[{{delta,wgt,-2 wgt/delta,D[-Sqrt[1+h^2],{h,2}] /. h -> 0}},
 TableHeadings -> {None,{"Gap","Weight","Retarded limit","Curvature"}}]
\end{wl}
The last column is the second derivative of the exact lower eigenvalue
of $H+hD$. Its agreement with the retarded limit uses a nondegenerate
ground state and a probe independent of field. A thermal mixture
also changes its populations, as the exercise below illustrates.

\section{Quasiparticle parameters and available phase space}
Write $G^{-1}=z-\epsilon_k+\mu-\Sigma(k,z)$. For a normal Fermi
liquid at a regular point of its Fermi surface, the pole strength is
$Z=[1-\partial_\omega\operatorname{Re}\Sigma(k_F,0)]^{-1}$.
The Fermi velocity is $Z[\nabla_k\epsilon_k+
\nabla_k\operatorname{Re}\Sigma(k,0)]_{k_F}$ in units $\hbar=1$.
Only for a momentum-independent self-energy can the entire mass
enhancement of an isotropic parabolic band be read off from $1/Z$.
At first order a static Hartree shift changes the Fermi energy but
gives no decay. At second order a particle can excite a particle--hole
pair; thermal factors restrict all three final quasiparticles to a
shell about the Fermi surface. Two independent low-energy integrations
remain after energy conservation. With nonsingular phase space in
three dimensions the retarded imaginary part is negative and of order
$-\big[\omega^2+(\pi T)^2\big]$. In two dimensions logarithms can
modify this law. Broadening the discrete lines of a small cluster is
not a measurement of that lifetime.
\begin{exercise}[Fixed frequency or static field?]\label{ex:response}
For a single free level of energy $\epsilon$ at temperature $T>0$,
find the retarded density commutator and the isothermal derivative
$\partial\langle n\rangle/\partial\epsilon$ for one spin. Explain why
the two answers differ.
\end{exercise}

\part{Function reference}
\chapter{Operators, algorithms and interfaces}
Entries are grouped by subject. Each gives the usage and details of
the corresponding help page and an evaluated example. Short operator
identities and selected matrices are printed with their results. See the alphabetical
index for symbols used in more than one subject.
\input{generated/reference}
\appendix
\chapter{Conventions at a glance}
\section*{Levels and indices}
\begin{tabular}{@{}ll@{}}\toprule
Input & Meaning\\\midrule
\fn{c[i]} & a spinful level, without operator-type or spin index\\
\fn{c[CR,i,UP]} & creation of its up electron\\
\fn{c[AN,i,DO]} & annihilation of its down electron\\
\fn{number[c[i]]} & sum over spin of the occupations\\
\fn{hubbard[c[i]]} & $n_{i\uparrow}n_{i\downarrow}$ without $U$\\
\fn{spinplus[c[i]]} & $c^\dagger_{i\uparrow}c_{i\downarrow}$\\
\fn{S[i,4]}, \fn{S[i,5]} & raising and lowering of a declared spin\\
\bottomrule\end{tabular}
\par\medskip
The internal type indices are \fn{CR}=0 and \fn{AN}=1; spin-$1/2$
labels are \fn{UP}=1, \fn{DO}=0. Do not infer bit order by sorting
these integers: consult the current \fn{BASIS}. A bare boson normally
needs no spin index. For spinless fermions use the mode convention of
the constructor in question; the generated explicit operator may
carry a trailing zero.

\section*{Order, charge and conjugation}
Normal ordering is a rule for reducing an operator, not an expectation
value. \fn{conj} reverses the order of factors and changes their type.
Complex coefficients must be conjugated too. The empty state has
$Q=-L$ for $L$ spin-degenerate levels; half filling has $Q=0$.
Highest-weight spin bases represent each complete spin multiplet once.
Reflection parity acts on states including the fermionic permutation
sign. Fermion parity $(-1)^N$ and spatial-reflection parity need not
agree. Energy gaps must specify which quantum numbers are held fixed.

\section*{Bloch, Nambu and retarded conventions}
An annihilator in real space has Fourier phase $+ikR$ and normalization
$N_s^{-1/2}$. Ordered orbitals within a cell fix the rows of $h(k)$.
Including their positions changes the phases of off-diagonal elements,
but neither the eigenvalues nor gauge-invariant observables. In the
particle-first Nambu spinor the hole block uses $-h^T(-k)$ for fermions;
the prefactor is $1/2$ when summing the full doubled basis.

For $\operatorname{Im}z>0$ the transform $e^{izt}$ of the retarded
kernel gives poles at $E_n-E_m$, with denominator $z-(E_n-E_m)$.
Fermionic diagonal spectral density integrates to one per canonical
mode including both signs of frequency. Use $\mathcal L A=[A,H]$;
exchanging the order of the commutator reflects all pole frequencies.

\chapter{Extending the operator algebra}
\setwlchapter{extensions}
\section{Normal order as a reduction rule}
To add an algebra, specify which pairs are out of order and what
replaces such a pair. For odd $a,b$ the replacement is $ab=-ba+\{a,b\}$.
For bosons use $ab=ba+[a,b]$. A scalar overlap on the right lowers
the operator degree and can therefore be factored out of \fn{nc}.
A new commutator rule without an ordering that eventually stops
interchanging pairs can instead give infinite recursion. In particular
do not define both directions of a swap to call each other.

\fn{ordering} sets the reference convention per family;
\fn{snegOrderedQ} is the predicate for canonical order.
\fn{acmt} and \fn{cmt} supply the elementary brackets, whereas
\fn{antikomutator} and \fn{komutator} construct brackets of arbitrary
products from the algebra. Test a rule in the declared operator family
first, preferably as an upvalue of its head; a broad downvalue on a
public symbol can change every subsequent calculation in that kernel.
\casestudy{custom}{A continuum field with a distributional overlap}

\section{Scalars and linear maps}
The predicate \fn{isnumericQ} is about commuting coefficients in
SNEG, not just numbers evaluable to floating point. A real dispersion
relation, a mesh weight and a distribution such as a Dirac delta
may all be scalar factors. Declaring an operator-valued function a
scalar would allow it to be commuted through its conjugate and
silently lose a physical term. Grassmann numbers require their own
exchange signs. When defining a linear map $F$, distinguish
$F(zA)=zF(A)$ from the conjugate-linear first slot of an inner product.

\section{Compatibility and tests}
User models should call the public symbols in \texttt{Sneg`}.
Internal helpers in \texttt{Sneg`Private`} are implementation details.
To qualify an extension, check the single-mode bracket, a three-factor
product that can be reduced at either adjacent pair, the adjoint of
both results, and a cross-family product. Then specialize to a finite
overlap or finite cluster with an independent matrix realization.
For example a positive Hermitian overlap matrix $G$ can be realized by
$d_i=\sum_j L_{ij}c_j$, $LL^\dagger=G$, using canonical fermions.
This checks conjugation and the ordering of the indices as well as signs.

Older interfaces can occur in existing research notebooks.
\fn{ZeroMatrix} is an old matrix-utility name; construct a square
zero matrix with \fn{ConstantArray[0,\{n,n\}]}. For package output in
a paper use the public \fn{snegtex}; retain \fn{InputForm} alongside
formatted output if another program must read index values. Numerical
comparisons between SNEG and NRG Ljubljana must use the same impurity
and chain-site order, parity convention and zero of energy.

\section{Interfaces used by existing models}
The following names may appear in a model file or in an extension to
the algebra. When calling a basis or contraction routine directly,
match the state representation and the normalization it expects.
\input{generated/compatibility}

\chapter{Solutions to the exercises}
\section*{Tutorial}
\setwlchapter{tutorial-solutions}
\begin{solution}{ex:first}
At $\epsilon=-U/3$ the empty, spin-doublet, and double energies
are $0,-U/3,U/3$. Starting from the doublet, removing costs $U/3$,
adding costs $2U/3$. Both excitations change charge and therefore
disappear in the restricted calculation. The spin degeneracy costs
no energy but does not supply a one-electron transition.
\begin{wl}
snegfermionoperators[d];
ha = hubbard[d[]] - number[d[]]/3;
en = #[[2,1,1]] & /@ makematricesbzvc[ha, qsbasisvc[{d[]}]];
TableForm[{{First[en]-en[[2]], Last[en]-en[[2]]}},
  TableHeadings -> {None, {"Removal / U", "Addition / U"}}]
\end{wl}
\end{solution}
\begin{solution}{ex:majorana}
The mixed terms cancel in $f^2$ and the equal squares enter with
opposite signs. In the anticommutator they add. In the number there is
an interference term $i\gamma_1\gamma_2$, which is Hermitian because
the two Majoranas anticommute. Direct reduction gives the results:
\begin{wl}
snegmajoranaoperators[gm];
f = (gm[1] + I gm[2])/Sqrt[2];
Expand[{nc[f,f], antikomutator[f,conj[f]], nc[conj[f],f]}]
\end{wl}
\end{solution}
\begin{solution}{ex:bond}
The molecule has one off-diagonal element $-t$, hence energies $\pm t$.
On the two-site ring the two translates contribute to the same matrix
element, giving $\pm2t$. For $L>2$ the two neighbours of a site are
distinct. Halving the hopping would halve the physical bandwidth.
\begin{wl}
snegfermionoperators[c];
bm = Select[qszbasisvc[{c[1],c[2]}], First[#] == {-1,1/2} &];
hy = -hop[c[1],c[2]];
Sort[Eigenvalues[#[[1,2]]]] & /@
  (makematricesbzvc[#,bm] & /@ {hy,2 hy}) // MatrixForm
\end{wl}
\end{solution}
\begin{solution}{ex:grid}
In the inner product the twist factors cancel pointwise. For column
indices $m,m'$ the remaining sum is $L^{-1}\sum_{n=0}^{L-1}
e^{2\pi i(m-m')n/L}=\delta_{mm'}$. The inverse is the adjoint, not a
second forward transform:
\begin{wl}
ff = Table[Exp[I n (2 Pi j + Pi/2)/3]/Sqrt[3], {n,0,2}, {j,0,2}];
FullSimplify[ConjugateTranspose[ff].ff] // MatrixForm
\end{wl}
\end{solution}
\begin{solution}{ex:multiplet}
With $x=e^{-J/T}$ the triplet carries probability $3x/(1+3x)$.
Thus $\langle\vS_1\cdot\vS_2\rangle=3(x-1)/[4(1+3x)]$, tending
to $-3/4$ at zero temperature and zero at infinite temperature.
Weighting the two energies equally would give a nonzero infinite-temperature
result. As a check evaluate the trace of the Boltzmann-weighted spin
matrix for $J/T=1$:
\begin{wl}
snegspinoperators[{s,1/2}];
eb = Flatten[Eigenvalues[Normal[Last[#]]] & /@
  spinclustermatrices[spinspin[s[1],s[2]], {s[1],s[2]}]];
FullSimplify[(eb.Exp[-eb])/Total[Exp[-eb]] == 3 (1-E)/(4 (E+3))]
\end{wl}
\end{solution}
\begin{solution}{ex:wick}
In the fixed-charge ensemble each spin has probability $1/2$, but
the two spins never occur together. Independent half-filled modes
have joint probability $1/4$. The difference is the connected
two-particle correlator. It can be checked by diagonal matrix elements
in the four Fock states:
\begin{wl}
bb = qbasisvc[{d[]}];
ddm = makematricesbzvc[hubbard[d[]],bb];
TableForm[{First[First[#]],Tr[Last[#]]/Length[Last[#]]} & /@ ddm,
  TableHeadings -> {None, {"Charge Q","Mean doublons"}}]
\end{wl}
The $Q=0$ row gives zero. Including the other sectors with their full
dimensions gives $(0+0+0+1)/4$.
\end{solution}
\begin{solution}{ex:trace}
Exactly one of the four local states is doubly occupied. At each site
there are $4^{L-1}$ such configurations; the trace is $U L4^{L-1}$.
Hopping has no diagonal element. $\mathrm{Tr}(H_t^2)$ sums squared
hopping matrix elements and reveals a missing bond, though not its
sign. For two spinful sites of unit hopping:
\begin{wl}
bb = qbasisvc[{c[1],c[2]}];
mm = Last /@ makematricesbzvc[hy,bb];
{Total[Tr /@ mm], Total[Tr[#.#] & /@ mm]}
\end{wl}
\end{solution}

\setwlchapter{correlated-solutions}

\section*{Correlated electrons}
The checks below form a fresh session. In each case the calculation uses
a different representation or route from the analytical argument.

\begin{solution}{ex:defects}
At fixed $N=L$, a state with no doublons has one electron on every
site. The spins can be chosen independently, giving $2^L$ ground
states. A state at energy $U$ has precisely one doublon and one hole.
Choose the doublon site in $L$ ways and the hole site in $L-1$ ways.
There is no spin choice on either of these two sites, giving
$L(L-1)2^{L-2}$ states. For four sites the two counts are 16 and 48.
The other six half-filled states have two doublons and two holes.
Their inclusion gives $\binom84=70$ as required. With $U$ as the
energy unit the full half-filled matrix is diagonal:
\begin{wl}
Needs["sneg`"];
snegfermionoperators[c, d];
sites4 = c /@ Range[4];
b4 = Select[qbasisvc[sites4], First[#] == {0} &];
atomic4 = First[makematricesbzvc[Total[hubbard[sites4]], b4]][[2]];
TableForm[Sort[Tally[Eigenvalues[atomic4]]],
  TableHeadings -> {None, {"Energy / U", "States"}}]
\end{wl}
\end{solution}

\begin{solution}{ex:detuned}
The singly occupied configurations have zero site energy. The two
intermediate doublons cost $U-\delta$ and $U+\delta$. Each is reached
from the singlet with squared amplitude $2t^2$, while the triplet
cannot hop. Consequently
\[
 J(\delta)=2t^2\left(\frac1{U-\delta}+\frac1{U+\delta}\right)
 =\frac{4t^2 U}{U^2-\delta^2}.
\]
Positive excitation costs require $|\delta|<U$, and we need
$|t|\ll U-|\delta|$. The divergence as $|\delta|\to U$ signals
that a charge state must be retained in the low-energy subspace.
It is not a physical divergence of the exchange energy. Verify at
finite detuning by diagonalizing the full singlet matrix; the two
answers are printed in the same energy units:
\begin{wl}
snegrealconstants[t, U, delta];
sites2 = {c[1], c[2]};
hd = U Total[hubbard[sites2]] - t hop[c[1], c[2]];
biased = hd + delta (number[c[2]] - number[c[1]])/2;
bs = Select[qsbasisvc[sites2], First[#] == {0, 0} &];
hs = First[makematricesbzvc[biased, bs]][[2]];
par = {U -> 10, delta -> 6, t -> 1/10};
j2 = 4 t^2 U/(U^2 - delta^2);
N[{-Min[Eigenvalues[hs /. par]], j2 /. par}, 6]
\end{wl}
The triplet energy is exactly zero even at finite detuning. The
difference between the two printed splittings is the contribution of
fourth and higher even orders.
\end{solution}

\begin{solution}{ex:gaps}
One electron occupies a bonding orbital of energy $-t$. A three-electron
state contains exactly one doublon, so its lowest energy is $U-t$.
The half-filled ground-state energy is $E_-$ of Eq.~\eqref{eq:dimer}.
Substitution in Eq.~\eqref{eq:chargegap} gives
\[
 \Delta_c=\sqrt{U^2+16t^2}-2t,\qquad
 \Delta_s=(\sqrt{U^2+16t^2}-U)/2.
\]
For $U=0$ both are $2t$: the bonding level forms a closed shell, and
the antibonding level lies $2t$ higher. At $U\gg t$,
$\Delta_c=U-2t+8t^2/U+\cdots$, while $\Delta_s=4t^2/U+\cdots$.
As a direct check, set $(t,U)=(1,4)$, use charge \emph{only}, and
take the lowest eigenvalue separately in each particle-number sector:
\begin{wl}
bc = qbasisvc[sites2];
mc = makematricesbzvc[hd /. {t -> 1, U -> 4}, bc];
e[q_] := Min[Eigenvalues[First[Select[mc, First[#] == {q} &]][[2]]]];
Simplify[e[1] + e[-1] - 2 e[0] == Sqrt[32] - 2]
\end{wl}
\end{solution}

\begin{solution}{ex:moments}
Expand $G_\sigma$ in inverse powers of $z$:
\[
 G_\sigma(z)=\sum_{\ell\ge0}M_\ell/z^{\ell+1}.
\]
Equation~\eqref{eq:atomicG} then gives
$M_\ell=(1-w)\epsilon^\ell+w(\epsilon+U)^\ell$. Thus
\[
 M_0=1,\quad M_1=\epsilon+Uw,\quad
 M_2=\epsilon^2+(2\epsilon U+U^2)w.
\]
A single pole at the Hartree energy reproduces $M_0$ and $M_1$,
but has $M_2=M_1^2$. The exact variance $M_2-M_1^2=U^2w(1-w)$
vanishes only when the peaks coincide or one carries zero weight.
For an independent computation \fn{spectralmoments} evaluates
$M_\ell=\langle\{ \mathcal{L}^\ell d_\uparrow,d^\dagger_\uparrow\}\rangle$
with $\mathcal{L}A=[A,H]$:
\begin{wl}
snegrealconstants[\[Epsilon]];
ha = \[Epsilon] number[d[]] + U hubbard[d[]];
mom = spectralmoments[d[AN, UP], d[CR, UP], ha, 2] /.
  avg[number[d[], DO]] -> w;
Expand[mom - {1, \[Epsilon] + U w, \[Epsilon]^2 + (2 \[Epsilon] U + U^2) w}]
\end{wl}
\end{solution}

\section*{Further applications}
\setwlchapter{application-solutions}
\begin{solution}{ex:harper}
The antisymmetric combination of sites 1 and 2 has energy 2 in
units of $t$. In the remaining subspace the trace is $-2$ and the
product of eigenvalues is $-2$, giving $-1\pm\sqrt3$. The sum of
all energies is zero, as is the sum of the diagonal potential in a
magnetic cell. The sum of their squares is $6+6=12$: six from the
on-site potentials and six from oriented links.
\begin{wl}
hm = {{-2,-1,-1},{-1,1,-1},{-1,-1,1}};
{Eigenvalues[hm], Tr[hm], Tr[hm.hm]}
\end{wl}
\end{solution}
\begin{solution}{ex:jt}
Minimize the lower-surface energy in the nonnegative radius $r$.
Its stationary point is $|g|/K$ with energy $-g^2/(2K)$ measured
from the undistorted crossing. For $K=M\omega^2$ the coefficient
of $b+b^\dagger$ is $g/\sqrt{2M\omega}$; the square divided by
$\omega$ is precisely the relaxation energy. Diagonalize the electronic
part on a circle with the same radius in two different directions:
\begin{wl}
el[x_,y_] := x PauliZ + y PauliX;
Sort[Eigenvalues[el[3,4]]] == Sort[Eigenvalues[el[0,5]]]
\end{wl}
\end{solution}
\begin{solution}{ex:andreev}
At particle--hole symmetry both diagonal even energies vanish.
Pair transfer splits this subspace into levels $-\Gamma,\Gamma$,
while each odd state is at $-U/2$. The doublet wins for $U>2\Gamma$;
at equality it is degenerate with the lower singlet. Test the crossing
at $\Gamma=1$ without imposing charge symmetry:
\begin{wl}
snegfermionoperators[d];
hp = 2 hubbard[d[]] - number[d[]] - bcs[d[]];
mm = makematricesbzvc[hp, nonebasisvc[{d[]}]];
Sort[Eigenvalues[First[mm][[2]]]]
\end{wl}
\end{solution}
\begin{solution}{ex:triangle}
Addition of three spins $1/2$ gives one $S=3/2$ quartet and two
$S=1/2$ doublets. Their energies are $3J/4$ and $-3J/4$, respectively.
Hence $Z=8\cosh(3J/4T)$ and the residual entropy is $\ln4$ in units
of $k_{\mathrm B}$. No choice of quantization axis changes that count:
\begin{wl}
snegspinoperators[{s,1/2}];
tri = Total[spinspin @@@ Subsets[s /@ Range[3], {2}]];
ev = Flatten[Eigenvalues[Normal[Last[#]]] & /@
  spinclustermatrices[tri,{s -> Range[3]}]];
TableForm[Sort[Tally[ev]], TableHeadings -> {None, {"Energy / J","States"}}]
\end{wl}
\end{solution}
\begin{solution}{ex:resonance}
Here $G_{dd}(z)=z/(z^2-V^2)$: each of the bonding and antibonding
orbitals has impurity probability $1/2$. The two half-weight poles
merge to the original unit-weight pole as $V\to0$. This limit is
not obtained by discarding one of the two levels.
\begin{wl}
snegfermionoperators[c];
snegrealconstants[V];
Apart[greensfunction[d[AN,UP],d[CR,UP],V hop[c[],d[]],z],z]
\end{wl}
\end{solution}
\begin{solution}{ex:response}
The occupation commutes with the free Hamiltonian, so its retarded
commutator vanishes identically. At thermal equilibrium
$\partial_\epsilon f=-f(1-f)/T$; changing the level energy changes
the Gibbs weights. For an isolated system a commuting probe cannot
redistribute them dynamically. At $T=1,\epsilon=0$ check the derivative
by a centred difference:
\begin{wl}
occ[e_] := thermalaverage[number[d[],UP], e number[d[]], {d}, 1];
{komutator[number[d[],UP],number[d[]]],
 Round[(occ[1/1000]-occ[-1/1000])/(2/1000),1/10000]}
\end{wl}
\end{solution}

\printbibliography[heading=bibintoc]
\addchap{Index of functions}
\input{generated/symbol-index}
\end{document}
